% Môn: Vật lý
% Lớp: 10
% Chương: Chương V. Động lượng
% Bài: L10C5 Bài 28. Động lượng
% Số câu: 162
% App đọc file này trực tiếp — sửa rồi Commit trên GitHub, bấm Đồng bộ GitHub .tex

% L10C5 Bài 28. Động lượng

% ===== Câu 1 =====
% ID: AIQ_L10C5_FULL_0001
% Mức: NB
\dangbt{Nhận biết động lượng và các đặc trưng của vectơ động lượng}
\begin{ex}
% ID: AIQ_L10C5_FULL_0001
% Mức: NB
Một vật khối lượng $1\ \mathrm{kg}$ chuyển động theo chiều dương với tốc độ $2\ \mathrm{m/s}$. Độ lớn động lượng bằng bao nhiêu?
\choice
{\True $2\ \mathrm{kg\,m/s}$}
{$3\ \mathrm{kg\,m/s}$}
{$1\ \mathrm{kg\,m/s}$}
{$20\ \mathrm{kg\,m/s}$}
\loigiai{
Áp dụng $p=mv=1\cdot2$. Suy ra kết quả $2\ \mathrm{kg\,m/s}$.
}
\end{ex}

% ===== Câu 2 =====
% ID: AIQ_L10C5_FULL_0002
% Mức: NB
\begin{ex}
% ID: AIQ_L10C5_FULL_0002
% Mức: NB
Một vật khối lượng $1,5\ \mathrm{kg}$ chuyển động theo chiều dương với tốc độ $3\ \mathrm{m/s}$. Độ lớn động lượng bằng bao nhiêu?
\choice
{Không đủ dữ kiện để có giá trị này}
{\True $4,5\ \mathrm{kg\,m/s}$}
{$1,5\ \mathrm{kg\,m/s}$}
{$45\ \mathrm{kg\,m/s}$}
\loigiai{
Áp dụng $p=mv=1,5\cdot3$. Suy ra kết quả $4,5\ \mathrm{kg\,m/s}$.
}
\end{ex}

% ===== Câu 3 =====
% ID: AIQ_L10C5_FULL_0003
% Mức: TH
\begin{ex}
% ID: AIQ_L10C5_FULL_0003
% Mức: TH
Một vật khối lượng $2\ \mathrm{kg}$ chuyển động theo chiều dương với tốc độ $4\ \mathrm{m/s}$. Độ lớn động lượng bằng bao nhiêu?
\choice
{$6\ \mathrm{kg\,m/s}$}
{$2\ \mathrm{kg\,m/s}$}
{\True $8\ \mathrm{kg\,m/s}$}
{$80\ \mathrm{kg\,m/s}$}
\loigiai{
Áp dụng $p=mv=2\cdot4$. Suy ra kết quả $8\ \mathrm{kg\,m/s}$.
}
\end{ex}

% ===== Câu 4 =====
% ID: AIQ_L10C5_FULL_0004
% Mức: TH
\begin{ex}
% ID: AIQ_L10C5_FULL_0004
% Mức: TH
Một vật khối lượng $2,5\ \mathrm{kg}$ chuyển động theo chiều dương với tốc độ $5\ \mathrm{m/s}$. Độ lớn động lượng bằng bao nhiêu?
\choice
{$7,5\ \mathrm{kg\,m/s}$}
{$2,5\ \mathrm{kg\,m/s}$}
{$125\ \mathrm{kg\,m/s}$}
{\True $12,5\ \mathrm{kg\,m/s}$}
\loigiai{
Áp dụng $p=mv=2,5\cdot5$. Suy ra kết quả $12,5\ \mathrm{kg\,m/s}$.
}
\end{ex}

% ===== Câu 5 =====
% ID: AIQ_L10C5_FULL_0005
% Mức: VD
\begin{ex}
% ID: AIQ_L10C5_FULL_0005
% Mức: VD
Một vật khối lượng $3\ \mathrm{kg}$ chuyển động theo chiều dương với tốc độ $6\ \mathrm{m/s}$. Độ lớn động lượng bằng bao nhiêu?
\choice
{\True $18\ \mathrm{kg\,m/s}$}
{$9\ \mathrm{kg\,m/s}$}
{$3\ \mathrm{kg\,m/s}$}
{$180\ \mathrm{kg\,m/s}$}
\loigiai{
Áp dụng $p=mv=3\cdot6$. Suy ra kết quả $18\ \mathrm{kg\,m/s}$.
}
\end{ex}

% ===== Câu 6 =====
% ID: AIQ_L10C5_FULL_0006
% Mức: VD
\begin{ex}
% ID: AIQ_L10C5_FULL_0006
% Mức: VD
Một vật khối lượng $1\ \mathrm{kg}$ chuyển động theo chiều dương với tốc độ $7\ \mathrm{m/s}$. Độ lớn động lượng bằng bao nhiêu?
\choice
{$8\ \mathrm{kg\,m/s}$}
{\True $7\ \mathrm{kg\,m/s}$}
{$6\ \mathrm{kg\,m/s}$}
{$70\ \mathrm{kg\,m/s}$}
\loigiai{
Áp dụng $p=mv=1\cdot7$. Suy ra kết quả $7\ \mathrm{kg\,m/s}$.
}
\end{ex}

% ===== Câu 7 =====
% ID: AIQ_L10C5_FULL_0007
% Mức: VD
\begin{ex}
% ID: AIQ_L10C5_FULL_0007
% Mức: VD
Một vật khối lượng $1,5\ \mathrm{kg}$ chuyển động theo chiều dương với tốc độ $8\ \mathrm{m/s}$. Độ lớn động lượng bằng bao nhiêu?
\choice
{$9,5\ \mathrm{kg\,m/s}$}
{$6,5\ \mathrm{kg\,m/s}$}
{\True $12\ \mathrm{kg\,m/s}$}
{$120\ \mathrm{kg\,m/s}$}
\loigiai{
Áp dụng $p=mv=1,5\cdot8$. Suy ra kết quả $12\ \mathrm{kg\,m/s}$.
}
\end{ex}

% ===== Câu 8 =====
% ID: AIQ_L10C5_FULL_0008
% Mức: VDC
\begin{ex}
% ID: AIQ_L10C5_FULL_0008
% Mức: VDC
Một vật khối lượng $2\ \mathrm{kg}$ chuyển động theo chiều dương với tốc độ $9\ \mathrm{m/s}$. Độ lớn động lượng bằng bao nhiêu?
\choice
{$11\ \mathrm{kg\,m/s}$}
{$7\ \mathrm{kg\,m/s}$}
{$180\ \mathrm{kg\,m/s}$}
{\True $18\ \mathrm{kg\,m/s}$}
\loigiai{
Áp dụng $p=mv=2\cdot9$. Suy ra kết quả $18\ \mathrm{kg\,m/s}$.
}
\end{ex}

% ===== Câu 9 =====
% ID: AIQ_L10C5_FULL_0009
% Mức: VDC
\begin{ex}
% ID: AIQ_L10C5_FULL_0009
% Mức: VDC
Một vật khối lượng $2,5\ \mathrm{kg}$ chuyển động theo chiều dương với tốc độ $10\ \mathrm{m/s}$. Độ lớn động lượng bằng bao nhiêu?
\choice
{\True $25\ \mathrm{kg\,m/s}$}
{$12,5\ \mathrm{kg\,m/s}$}
{$7,5\ \mathrm{kg\,m/s}$}
{$250\ \mathrm{kg\,m/s}$}
\loigiai{
Áp dụng $p=mv=2,5\cdot10$. Suy ra kết quả $25\ \mathrm{kg\,m/s}$.
}
\end{ex}

% ===== Câu 10 =====
% ID: AIQ_L10C5_FULL_0010
% Mức: VDC
\begin{ex}
% ID: AIQ_L10C5_FULL_0010
% Mức: VDC
Một vật khối lượng $3\ \mathrm{kg}$ chuyển động theo chiều dương với tốc độ $11\ \mathrm{m/s}$. Độ lớn động lượng bằng bao nhiêu?
\choice
{$14\ \mathrm{kg\,m/s}$}
{\True $33\ \mathrm{kg\,m/s}$}
{$8\ \mathrm{kg\,m/s}$}
{$330\ \mathrm{kg\,m/s}$}
\loigiai{
Áp dụng $p=mv=3\cdot11$. Suy ra kết quả $33\ \mathrm{kg\,m/s}$.
}
\end{ex}

% ===== Câu 11 =====
% ID: AIQ_L10C5_FULL_0011
% Mức: NB
\begin{ex}
% ID: AIQ_L10C5_FULL_0011
% Mức: NB
Xét các phát biểu thuộc dạng “Nhận biết động lượng và các đặc trưng của vectơ động lượng” ở tình huống số 1.
\choiceTF
{\True Động lượng là đại lượng vectơ.}
{\True Vectơ động lượng cùng hướng với vectơ vận tốc.}
{Đơn vị động lượng là N.}
{\True Động lượng phụ thuộc hệ quy chiếu.}
\loigiai{
\begin{itemchoice}
\item (Đúng) Động lượng là đại lượng vectơ. (Vì): Phù hợp với định nghĩa hoặc hệ thức vật lí. b) Đúng: Phù hợp với định nghĩa hoặc hệ thức vật lí. c) Sai: Không phù hợp với định nghĩa hoặc hệ thức vật lí. d) Đúng: Phù hợp với định nghĩa hoặc hệ thức vật lí
\item (Đúng) Vectơ động lượng cùng hướng với vectơ vận tốc.
\item (Sai) Đơn vị động lượng là N.
\item (Đúng) Động lượng phụ thuộc hệ quy chiếu.
\end{itemchoice}
}
\end{ex}

% ===== Câu 12 =====
% ID: AIQ_L10C5_FULL_0012
% Mức: TH
\begin{ex}
% ID: AIQ_L10C5_FULL_0012
% Mức: TH
Xét các phát biểu thuộc dạng “Nhận biết động lượng và các đặc trưng của vectơ động lượng” ở tình huống số 2.
\choiceTF
{\True Động lượng là đại lượng vectơ.}
{\True Vectơ động lượng cùng hướng với vectơ vận tốc.}
{Đơn vị động lượng là N.}
{\True Động lượng phụ thuộc hệ quy chiếu.}
\loigiai{
\begin{itemchoice}
\item (Đúng) Động lượng là đại lượng vectơ. (Vì): Phù hợp với định nghĩa hoặc hệ thức vật lí. b) Đúng: Phù hợp với định nghĩa hoặc hệ thức vật lí. c) Sai: Không phù hợp với định nghĩa hoặc hệ thức vật lí. d) Đúng: Phù hợp với định nghĩa hoặc hệ thức vật lí
\item (Đúng) Vectơ động lượng cùng hướng với vectơ vận tốc.
\item (Sai) Đơn vị động lượng là N.
\item (Đúng) Động lượng phụ thuộc hệ quy chiếu.
\end{itemchoice}
}
\end{ex}

% ===== Câu 13 =====
% ID: AIQ_L10C5_FULL_0013
% Mức: TH
\begin{ex}
% ID: AIQ_L10C5_FULL_0013
% Mức: TH
Xét các phát biểu thuộc dạng “Nhận biết động lượng và các đặc trưng của vectơ động lượng” ở tình huống số 3.
\choiceTF
{\True Động lượng là đại lượng vectơ.}
{\True Vectơ động lượng cùng hướng với vectơ vận tốc.}
{Đơn vị động lượng là N.}
{\True Động lượng phụ thuộc hệ quy chiếu.}
\loigiai{
\begin{itemchoice}
\item (Đúng) Động lượng là đại lượng vectơ. (Vì): Phù hợp với định nghĩa hoặc hệ thức vật lí. b) Đúng: Phù hợp với định nghĩa hoặc hệ thức vật lí. c) Sai: Không phù hợp với định nghĩa hoặc hệ thức vật lí. d) Đúng: Phù hợp với định nghĩa hoặc hệ thức vật lí
\item (Đúng) Vectơ động lượng cùng hướng với vectơ vận tốc.
\item (Sai) Đơn vị động lượng là N.
\item (Đúng) Động lượng phụ thuộc hệ quy chiếu.
\end{itemchoice}
}
\end{ex}

% ===== Câu 14 =====
% ID: AIQ_L10C5_FULL_0014
% Mức: VD
\begin{ex}
% ID: AIQ_L10C5_FULL_0014
% Mức: VD
Xét các phát biểu thuộc dạng “Nhận biết động lượng và các đặc trưng của vectơ động lượng” ở tình huống số 4.
\choiceTF
{\True Động lượng là đại lượng vectơ.}
{\True Vectơ động lượng cùng hướng với vectơ vận tốc.}
{Đơn vị động lượng là N.}
{\True Động lượng phụ thuộc hệ quy chiếu.}
\loigiai{
\begin{itemchoice}
\item (Đúng) Động lượng là đại lượng vectơ. (Vì): Phù hợp với định nghĩa hoặc hệ thức vật lí. b) Đúng: Phù hợp với định nghĩa hoặc hệ thức vật lí. c) Sai: Không phù hợp với định nghĩa hoặc hệ thức vật lí. d) Đúng: Phù hợp với định nghĩa hoặc hệ thức vật lí
\item (Đúng) Vectơ động lượng cùng hướng với vectơ vận tốc.
\item (Sai) Đơn vị động lượng là N.
\item (Đúng) Động lượng phụ thuộc hệ quy chiếu.
\end{itemchoice}
}
\end{ex}

% ===== Câu 15 =====
% ID: AIQ_L10C5_FULL_0015
% Mức: VDC
\begin{ex}
% ID: AIQ_L10C5_FULL_0015
% Mức: VDC
Xét các phát biểu thuộc dạng “Nhận biết động lượng và các đặc trưng của vectơ động lượng” ở tình huống số 5.
\choiceTF
{\True Động lượng là đại lượng vectơ.}
{\True Vectơ động lượng cùng hướng với vectơ vận tốc.}
{Đơn vị động lượng là N.}
{\True Động lượng phụ thuộc hệ quy chiếu.}
\loigiai{
\begin{itemchoice}
\item (Đúng) Động lượng là đại lượng vectơ. (Vì): Phù hợp với định nghĩa hoặc hệ thức vật lí. b) Đúng: Phù hợp với định nghĩa hoặc hệ thức vật lí. c) Sai: Không phù hợp với định nghĩa hoặc hệ thức vật lí. d) Đúng: Phù hợp với định nghĩa hoặc hệ thức vật lí
\item (Đúng) Vectơ động lượng cùng hướng với vectơ vận tốc.
\item (Sai) Đơn vị động lượng là N.
\item (Đúng) Động lượng phụ thuộc hệ quy chiếu.
\end{itemchoice}
}
\end{ex}

% ===== Câu 16 =====
% ID: AIQ_L10C5_FULL_0016
% Mức: TH
\begin{ex}
% ID: AIQ_L10C5_FULL_0016
% Mức: TH
Một vật khối lượng $1\ \mathrm{kg}$ chuyển động theo chiều dương với tốc độ $12\ \mathrm{m/s}$. Độ lớn động lượng bằng bao nhiêu?
\shortans{12}
\loigiai{
Áp dụng $p=mv=1\cdot12$. Suy ra kết quả $12\ \mathrm{kg\,m/s}$. Đáp án trả lời ngắn không ghi đơn vị.
}
\end{ex}

% ===== Câu 17 =====
% ID: AIQ_L10C5_FULL_0017
% Mức: TH
\begin{ex}
% ID: AIQ_L10C5_FULL_0017
% Mức: TH
Một vật khối lượng $1,5\ \mathrm{kg}$ chuyển động theo chiều dương với tốc độ $13\ \mathrm{m/s}$. Độ lớn động lượng bằng bao nhiêu?
\shortans{19,5}
\loigiai{
Áp dụng $p=mv=1,5\cdot13$. Suy ra kết quả $19,5\ \mathrm{kg\,m/s}$. Đáp án trả lời ngắn không ghi đơn vị.
}
\end{ex}

% ===== Câu 18 =====
% ID: AIQ_L10C5_FULL_0018
% Mức: VD
\begin{ex}
% ID: AIQ_L10C5_FULL_0018
% Mức: VD
Một vật khối lượng $2\ \mathrm{kg}$ chuyển động theo chiều dương với tốc độ $14\ \mathrm{m/s}$. Độ lớn động lượng bằng bao nhiêu?
\shortans{28}
\loigiai{
Áp dụng $p=mv=2\cdot14$. Suy ra kết quả $28\ \mathrm{kg\,m/s}$. Đáp án trả lời ngắn không ghi đơn vị.
}
\end{ex}

% ===== Câu 19 =====
% ID: AIQ_L10C5_FULL_0019
% Mức: VD
\begin{ex}
% ID: AIQ_L10C5_FULL_0019
% Mức: VD
Một vật khối lượng $2,5\ \mathrm{kg}$ chuyển động theo chiều dương với tốc độ $15\ \mathrm{m/s}$. Độ lớn động lượng bằng bao nhiêu?
\shortans{37,5}
\loigiai{
Áp dụng $p=mv=2,5\cdot15$. Suy ra kết quả $37,5\ \mathrm{kg\,m/s}$. Đáp án trả lời ngắn không ghi đơn vị.
}
\end{ex}

% ===== Câu 20 =====
% ID: AIQ_L10C5_FULL_0020
% Mức: VDC
\begin{ex}
% ID: AIQ_L10C5_FULL_0020
% Mức: VDC
Một vật khối lượng $3\ \mathrm{kg}$ chuyển động theo chiều dương với tốc độ $16\ \mathrm{m/s}$. Độ lớn động lượng bằng bao nhiêu?
\shortans{48}
\loigiai{
Áp dụng $p=mv=3\cdot16$. Suy ra kết quả $48\ \mathrm{kg\,m/s}$. Đáp án trả lời ngắn không ghi đơn vị.
}
\end{ex}

% ===== Câu 21 =====
% ID: AIQ_L10C5_FULL_0021
% Mức: VDC
\begin{ex}
% ID: AIQ_L10C5_FULL_0021
% Mức: VDC
Một vật khối lượng $1\ \mathrm{kg}$ chuyển động theo chiều dương với tốc độ $17\ \mathrm{m/s}$. Độ lớn động lượng bằng bao nhiêu?
\shortans{17}
\loigiai{
Áp dụng $p=mv=1\cdot17$. Suy ra kết quả $17\ \mathrm{kg\,m/s}$. Đáp án trả lời ngắn không ghi đơn vị.
}
\end{ex}

% ===== Câu 22 =====
% ID: AIQ_L10C5_FULL_0022
% Mức: TH
\begin{bt}
% ID: AIQ_L10C5_FULL_0022
% Mức: TH
Trình bày cách giải: Một vật khối lượng $1,5\ \mathrm{kg}$ chuyển động theo chiều dương với tốc độ $18\ \mathrm{m/s}$. Độ lớn động lượng bằng bao nhiêu?
\loigiai{
Áp dụng $p=mv=1,5\cdot18$. Suy ra kết quả $27\ \mathrm{kg\,m/s}$.
}
\end{bt}

% ===== Câu 23 =====
% ID: AIQ_L10C5_FULL_0023
% Mức: TH
\begin{bt}
% ID: AIQ_L10C5_FULL_0023
% Mức: TH
Trình bày cách giải: Một vật khối lượng $2\ \mathrm{kg}$ chuyển động theo chiều dương với tốc độ $19\ \mathrm{m/s}$. Độ lớn động lượng bằng bao nhiêu?
\loigiai{
Áp dụng $p=mv=2\cdot19$. Suy ra kết quả $38\ \mathrm{kg\,m/s}$.
}
\end{bt}

% ===== Câu 24 =====
% ID: AIQ_L10C5_FULL_0024
% Mức: VD
\begin{bt}
% ID: AIQ_L10C5_FULL_0024
% Mức: VD
Trình bày cách giải: Một vật khối lượng $2,5\ \mathrm{kg}$ chuyển động theo chiều dương với tốc độ $20\ \mathrm{m/s}$. Độ lớn động lượng bằng bao nhiêu?
\loigiai{
Áp dụng $p=mv=2,5\cdot20$. Suy ra kết quả $50\ \mathrm{kg\,m/s}$.
}
\end{bt}

% ===== Câu 25 =====
% ID: AIQ_L10C5_FULL_0025
% Mức: VD
\begin{bt}
% ID: AIQ_L10C5_FULL_0025
% Mức: VD
Trình bày cách giải: Một vật khối lượng $3\ \mathrm{kg}$ chuyển động theo chiều dương với tốc độ $21\ \mathrm{m/s}$. Độ lớn động lượng bằng bao nhiêu?
\loigiai{
Áp dụng $p=mv=3\cdot21$. Suy ra kết quả $63\ \mathrm{kg\,m/s}$.
}
\end{bt}

% ===== Câu 26 =====
% ID: AIQ_L10C5_FULL_0026
% Mức: VDC
\begin{bt}
% ID: AIQ_L10C5_FULL_0026
% Mức: VDC
Trình bày cách giải: Một vật khối lượng $1\ \mathrm{kg}$ chuyển động theo chiều dương với tốc độ $22\ \mathrm{m/s}$. Độ lớn động lượng bằng bao nhiêu?
\loigiai{
Áp dụng $p=mv=1\cdot22$. Suy ra kết quả $22\ \mathrm{kg\,m/s}$.
}
\end{bt}

% ===== Câu 27 =====
% ID: AIQ_L10C5_FULL_0027
% Mức: VDC
\begin{bt}
% ID: AIQ_L10C5_FULL_0027
% Mức: VDC
Trình bày cách giải: Một vật khối lượng $1,5\ \mathrm{kg}$ chuyển động theo chiều dương với tốc độ $23\ \mathrm{m/s}$. Độ lớn động lượng bằng bao nhiêu?
\loigiai{
Áp dụng $p=mv=1,5\cdot23$. Suy ra kết quả $34,5\ \mathrm{kg\,m/s}$.
}
\end{bt}

% ===== Câu 28 =====
% ID: AIQ_L10C5_FULL_0028
% Mức: NB
\dangbt{Tính động lượng của một vật}
\begin{ex}
% ID: AIQ_L10C5_FULL_0028
% Mức: NB
Xe nhỏ khối lượng $2\ \mathrm{kg}$ chuyển động với tốc độ $3\ \mathrm{m/s}$. Tính động lượng của xe.
\choice
{\True $6\ \mathrm{kg\,m/s}$}
{$5\ \mathrm{kg\,m/s}$}
{$3\ \mathrm{kg\,m/s}$}
{$9\ \mathrm{kg\,m/s}$}
\loigiai{
$p=mv=2\cdot3$. Suy ra kết quả $6\ \mathrm{kg\,m/s}$.
}
\end{ex}

% ===== Câu 29 =====
% ID: AIQ_L10C5_FULL_0029
% Mức: NB
\begin{ex}
% ID: AIQ_L10C5_FULL_0029
% Mức: NB
Xe nhỏ khối lượng $3\ \mathrm{kg}$ chuyển động với tốc độ $4\ \mathrm{m/s}$. Tính động lượng của xe.
\choice
{$7\ \mathrm{kg\,m/s}$}
{\True $12\ \mathrm{kg\,m/s}$}
{$6\ \mathrm{kg\,m/s}$}
{$16\ \mathrm{kg\,m/s}$}
\loigiai{
$p=mv=3\cdot4$. Suy ra kết quả $12\ \mathrm{kg\,m/s}$.
}
\end{ex}

% ===== Câu 30 =====
% ID: AIQ_L10C5_FULL_0030
% Mức: TH
\begin{ex}
% ID: AIQ_L10C5_FULL_0030
% Mức: TH
Xe nhỏ khối lượng $4\ \mathrm{kg}$ chuyển động với tốc độ $5\ \mathrm{m/s}$. Tính động lượng của xe.
\choice
{$9\ \mathrm{kg\,m/s}$}
{$10\ \mathrm{kg\,m/s}$}
{\True $20\ \mathrm{kg\,m/s}$}
{$25\ \mathrm{kg\,m/s}$}
\loigiai{
$p=mv=4\cdot5$. Suy ra kết quả $20\ \mathrm{kg\,m/s}$.
}
\end{ex}

% ===== Câu 31 =====
% ID: AIQ_L10C5_FULL_0031
% Mức: TH
\begin{ex}
% ID: AIQ_L10C5_FULL_0031
% Mức: TH
Xe nhỏ khối lượng $5\ \mathrm{kg}$ chuyển động với tốc độ $6\ \mathrm{m/s}$. Tính động lượng của xe.
\choice
{$11\ \mathrm{kg\,m/s}$}
{$15\ \mathrm{kg\,m/s}$}
{$36\ \mathrm{kg\,m/s}$}
{\True $30\ \mathrm{kg\,m/s}$}
\loigiai{
$p=mv=5\cdot6$. Suy ra kết quả $30\ \mathrm{kg\,m/s}$.
}
\end{ex}

% ===== Câu 32 =====
% ID: AIQ_L10C5_FULL_0032
% Mức: VD
\begin{ex}
% ID: AIQ_L10C5_FULL_0032
% Mức: VD
Xe nhỏ khối lượng $2\ \mathrm{kg}$ chuyển động với tốc độ $7\ \mathrm{m/s}$. Tính động lượng của xe.
\choice
{\True $14\ \mathrm{kg\,m/s}$}
{$9\ \mathrm{kg\,m/s}$}
{$7\ \mathrm{kg\,m/s}$}
{$21\ \mathrm{kg\,m/s}$}
\loigiai{
$p=mv=2\cdot7$. Suy ra kết quả $14\ \mathrm{kg\,m/s}$.
}
\end{ex}

% ===== Câu 33 =====
% ID: AIQ_L10C5_FULL_0033
% Mức: VD
\begin{ex}
% ID: AIQ_L10C5_FULL_0033
% Mức: VD
Xe nhỏ khối lượng $3\ \mathrm{kg}$ chuyển động với tốc độ $8\ \mathrm{m/s}$. Tính động lượng của xe.
\choice
{$11\ \mathrm{kg\,m/s}$}
{\True $24\ \mathrm{kg\,m/s}$}
{$12\ \mathrm{kg\,m/s}$}
{$32\ \mathrm{kg\,m/s}$}
\loigiai{
$p=mv=3\cdot8$. Suy ra kết quả $24\ \mathrm{kg\,m/s}$.
}
\end{ex}

% ===== Câu 34 =====
% ID: AIQ_L10C5_FULL_0034
% Mức: VD
\begin{ex}
% ID: AIQ_L10C5_FULL_0034
% Mức: VD
Xe nhỏ khối lượng $4\ \mathrm{kg}$ chuyển động với tốc độ $9\ \mathrm{m/s}$. Tính động lượng của xe.
\choice
{$13\ \mathrm{kg\,m/s}$}
{$18\ \mathrm{kg\,m/s}$}
{\True $36\ \mathrm{kg\,m/s}$}
{$45\ \mathrm{kg\,m/s}$}
\loigiai{
$p=mv=4\cdot9$. Suy ra kết quả $36\ \mathrm{kg\,m/s}$.
}
\end{ex}

% ===== Câu 35 =====
% ID: AIQ_L10C5_FULL_0035
% Mức: VDC
\begin{ex}
% ID: AIQ_L10C5_FULL_0035
% Mức: VDC
Xe nhỏ khối lượng $5\ \mathrm{kg}$ chuyển động với tốc độ $10\ \mathrm{m/s}$. Tính động lượng của xe.
\choice
{$15\ \mathrm{kg\,m/s}$}
{$25\ \mathrm{kg\,m/s}$}
{$60\ \mathrm{kg\,m/s}$}
{\True $50\ \mathrm{kg\,m/s}$}
\loigiai{
$p=mv=5\cdot10$. Suy ra kết quả $50\ \mathrm{kg\,m/s}$.
}
\end{ex}

% ===== Câu 36 =====
% ID: AIQ_L10C5_FULL_0036
% Mức: VDC
\begin{ex}
% ID: AIQ_L10C5_FULL_0036
% Mức: VDC
Xe nhỏ khối lượng $2\ \mathrm{kg}$ chuyển động với tốc độ $11\ \mathrm{m/s}$. Tính động lượng của xe.
\choice
{\True $22\ \mathrm{kg\,m/s}$}
{$13\ \mathrm{kg\,m/s}$}
{$11\ \mathrm{kg\,m/s}$}
{$33\ \mathrm{kg\,m/s}$}
\loigiai{
$p=mv=2\cdot11$. Suy ra kết quả $22\ \mathrm{kg\,m/s}$.
}
\end{ex}

% ===== Câu 37 =====
% ID: AIQ_L10C5_FULL_0037
% Mức: VDC
\begin{ex}
% ID: AIQ_L10C5_FULL_0037
% Mức: VDC
Xe nhỏ khối lượng $3\ \mathrm{kg}$ chuyển động với tốc độ $12\ \mathrm{m/s}$. Tính động lượng của xe.
\choice
{$15\ \mathrm{kg\,m/s}$}
{\True $36\ \mathrm{kg\,m/s}$}
{$18\ \mathrm{kg\,m/s}$}
{$48\ \mathrm{kg\,m/s}$}
\loigiai{
$p=mv=3\cdot12$. Suy ra kết quả $36\ \mathrm{kg\,m/s}$.
}
\end{ex}

% ===== Câu 38 =====
% ID: AIQ_L10C5_FULL_0038
% Mức: NB
\begin{ex}
% ID: AIQ_L10C5_FULL_0038
% Mức: NB
Xét các phát biểu thuộc dạng “Tính động lượng của một vật” ở tình huống số 1.
\choiceTF
{\True Công thức độ lớn động lượng là $p=mv$.}
{Khối lượng tăng thì động lượng giảm khi vận tốc không đổi.}
{\True Động lượng có đơn vị kg m/s.}
{\True Vật đứng yên có động lượng bằng 0.}
\loigiai{
\begin{itemchoice}
\item (Đúng) Công thức độ lớn động lượng là $p=mv$. (Vì): Phù hợp với định nghĩa hoặc hệ thức vật lí. b) Sai: Không phù hợp với định nghĩa hoặc hệ thức vật lí. c) Đúng: Phù hợp với định nghĩa hoặc hệ thức vật lí. d) Đúng: Phù hợp với định nghĩa hoặc hệ thức vật lí
\item (Sai) Khối lượng tăng thì động lượng giảm khi vận tốc không đổi.
\item (Đúng) Động lượng có đơn vị kg m/s.
\item (Đúng) Vật đứng yên có động lượng bằng 0.
\end{itemchoice}
}
\end{ex}

% ===== Câu 39 =====
% ID: AIQ_L10C5_FULL_0039
% Mức: TH
\begin{ex}
% ID: AIQ_L10C5_FULL_0039
% Mức: TH
Xét các phát biểu thuộc dạng “Tính động lượng của một vật” ở tình huống số 2.
\choiceTF
{\True Công thức độ lớn động lượng là $p=mv$.}
{Khối lượng tăng thì động lượng giảm khi vận tốc không đổi.}
{\True Động lượng có đơn vị kg m/s.}
{\True Vật đứng yên có động lượng bằng 0.}
\loigiai{
\begin{itemchoice}
\item (Đúng) Công thức độ lớn động lượng là $p=mv$. (Vì): Phù hợp với định nghĩa hoặc hệ thức vật lí. b) Sai: Không phù hợp với định nghĩa hoặc hệ thức vật lí. c) Đúng: Phù hợp với định nghĩa hoặc hệ thức vật lí. d) Đúng: Phù hợp với định nghĩa hoặc hệ thức vật lí
\item (Sai) Khối lượng tăng thì động lượng giảm khi vận tốc không đổi.
\item (Đúng) Động lượng có đơn vị kg m/s.
\item (Đúng) Vật đứng yên có động lượng bằng 0.
\end{itemchoice}
}
\end{ex}

% ===== Câu 40 =====
% ID: AIQ_L10C5_FULL_0040
% Mức: TH
\begin{ex}
% ID: AIQ_L10C5_FULL_0040
% Mức: TH
Xét các phát biểu thuộc dạng “Tính động lượng của một vật” ở tình huống số 3.
\choiceTF
{\True Công thức độ lớn động lượng là $p=mv$.}
{Khối lượng tăng thì động lượng giảm khi vận tốc không đổi.}
{\True Động lượng có đơn vị kg m/s.}
{\True Vật đứng yên có động lượng bằng 0.}
\loigiai{
\begin{itemchoice}
\item (Đúng) Công thức độ lớn động lượng là $p=mv$. (Vì): Phù hợp với định nghĩa hoặc hệ thức vật lí. b) Sai: Không phù hợp với định nghĩa hoặc hệ thức vật lí. c) Đúng: Phù hợp với định nghĩa hoặc hệ thức vật lí. d) Đúng: Phù hợp với định nghĩa hoặc hệ thức vật lí
\item (Sai) Khối lượng tăng thì động lượng giảm khi vận tốc không đổi.
\item (Đúng) Động lượng có đơn vị kg m/s.
\item (Đúng) Vật đứng yên có động lượng bằng 0.
\end{itemchoice}
}
\end{ex}

% ===== Câu 41 =====
% ID: AIQ_L10C5_FULL_0041
% Mức: VD
\begin{ex}
% ID: AIQ_L10C5_FULL_0041
% Mức: VD
Xét các phát biểu thuộc dạng “Tính động lượng của một vật” ở tình huống số 4.
\choiceTF
{\True Công thức độ lớn động lượng là $p=mv$.}
{Khối lượng tăng thì động lượng giảm khi vận tốc không đổi.}
{\True Động lượng có đơn vị kg m/s.}
{\True Vật đứng yên có động lượng bằng 0.}
\loigiai{
\begin{itemchoice}
\item (Đúng) Công thức độ lớn động lượng là $p=mv$. (Vì): Phù hợp với định nghĩa hoặc hệ thức vật lí. b) Sai: Không phù hợp với định nghĩa hoặc hệ thức vật lí. c) Đúng: Phù hợp với định nghĩa hoặc hệ thức vật lí. d) Đúng: Phù hợp với định nghĩa hoặc hệ thức vật lí
\item (Sai) Khối lượng tăng thì động lượng giảm khi vận tốc không đổi.
\item (Đúng) Động lượng có đơn vị kg m/s.
\item (Đúng) Vật đứng yên có động lượng bằng 0.
\end{itemchoice}
}
\end{ex}

% ===== Câu 42 =====
% ID: AIQ_L10C5_FULL_0042
% Mức: VDC
\begin{ex}
% ID: AIQ_L10C5_FULL_0042
% Mức: VDC
Xét các phát biểu thuộc dạng “Tính động lượng của một vật” ở tình huống số 5.
\choiceTF
{\True Công thức độ lớn động lượng là $p=mv$.}
{Khối lượng tăng thì động lượng giảm khi vận tốc không đổi.}
{\True Động lượng có đơn vị kg m/s.}
{\True Vật đứng yên có động lượng bằng 0.}
\loigiai{
\begin{itemchoice}
\item (Đúng) Công thức độ lớn động lượng là $p=mv$. (Vì): Phù hợp với định nghĩa hoặc hệ thức vật lí. b) Sai: Không phù hợp với định nghĩa hoặc hệ thức vật lí. c) Đúng: Phù hợp với định nghĩa hoặc hệ thức vật lí. d) Đúng: Phù hợp với định nghĩa hoặc hệ thức vật lí
\item (Sai) Khối lượng tăng thì động lượng giảm khi vận tốc không đổi.
\item (Đúng) Động lượng có đơn vị kg m/s.
\item (Đúng) Vật đứng yên có động lượng bằng 0.
\end{itemchoice}
}
\end{ex}

% ===== Câu 43 =====
% ID: AIQ_L10C5_FULL_0043
% Mức: TH
\begin{ex}
% ID: AIQ_L10C5_FULL_0043
% Mức: TH
Xe nhỏ khối lượng $4\ \mathrm{kg}$ chuyển động với tốc độ $13\ \mathrm{m/s}$. Tính động lượng của xe.
\shortans{52}
\loigiai{
$p=mv=4\cdot13$. Suy ra kết quả $52\ \mathrm{kg\,m/s}$. Đáp án trả lời ngắn không ghi đơn vị.
}
\end{ex}

% ===== Câu 44 =====
% ID: AIQ_L10C5_FULL_0044
% Mức: TH
\begin{ex}
% ID: AIQ_L10C5_FULL_0044
% Mức: TH
Xe nhỏ khối lượng $5\ \mathrm{kg}$ chuyển động với tốc độ $14\ \mathrm{m/s}$. Tính động lượng của xe.
\shortans{70}
\loigiai{
$p=mv=5\cdot14$. Suy ra kết quả $70\ \mathrm{kg\,m/s}$. Đáp án trả lời ngắn không ghi đơn vị.
}
\end{ex}

% ===== Câu 45 =====
% ID: AIQ_L10C5_FULL_0045
% Mức: VD
\begin{ex}
% ID: AIQ_L10C5_FULL_0045
% Mức: VD
Xe nhỏ khối lượng $2\ \mathrm{kg}$ chuyển động với tốc độ $15\ \mathrm{m/s}$. Tính động lượng của xe.
\shortans{30}
\loigiai{
$p=mv=2\cdot15$. Suy ra kết quả $30\ \mathrm{kg\,m/s}$. Đáp án trả lời ngắn không ghi đơn vị.
}
\end{ex}

% ===== Câu 46 =====
% ID: AIQ_L10C5_FULL_0046
% Mức: VD
\begin{ex}
% ID: AIQ_L10C5_FULL_0046
% Mức: VD
Xe nhỏ khối lượng $3\ \mathrm{kg}$ chuyển động với tốc độ $16\ \mathrm{m/s}$. Tính động lượng của xe.
\shortans{48}
\loigiai{
$p=mv=3\cdot16$. Suy ra kết quả $48\ \mathrm{kg\,m/s}$. Đáp án trả lời ngắn không ghi đơn vị.
}
\end{ex}

% ===== Câu 47 =====
% ID: AIQ_L10C5_FULL_0047
% Mức: VDC
\begin{ex}
% ID: AIQ_L10C5_FULL_0047
% Mức: VDC
Xe nhỏ khối lượng $4\ \mathrm{kg}$ chuyển động với tốc độ $17\ \mathrm{m/s}$. Tính động lượng của xe.
\shortans{68}
\loigiai{
$p=mv=4\cdot17$. Suy ra kết quả $68\ \mathrm{kg\,m/s}$. Đáp án trả lời ngắn không ghi đơn vị.
}
\end{ex}

% ===== Câu 48 =====
% ID: AIQ_L10C5_FULL_0048
% Mức: VDC
\begin{ex}
% ID: AIQ_L10C5_FULL_0048
% Mức: VDC
Xe nhỏ khối lượng $5\ \mathrm{kg}$ chuyển động với tốc độ $18\ \mathrm{m/s}$. Tính động lượng của xe.
\shortans{90}
\loigiai{
$p=mv=5\cdot18$. Suy ra kết quả $90\ \mathrm{kg\,m/s}$. Đáp án trả lời ngắn không ghi đơn vị.
}
\end{ex}

% ===== Câu 49 =====
% ID: AIQ_L10C5_FULL_0049
% Mức: TH
\begin{bt}
% ID: AIQ_L10C5_FULL_0049
% Mức: TH
Trình bày cách giải: Xe nhỏ khối lượng $2\ \mathrm{kg}$ chuyển động với tốc độ $19\ \mathrm{m/s}$. Tính động lượng của xe.
\loigiai{
$p=mv=2\cdot19$. Suy ra kết quả $38\ \mathrm{kg\,m/s}$.
}
\end{bt}

% ===== Câu 50 =====
% ID: AIQ_L10C5_FULL_0050
% Mức: TH
\begin{bt}
% ID: AIQ_L10C5_FULL_0050
% Mức: TH
Trình bày cách giải: Xe nhỏ khối lượng $3\ \mathrm{kg}$ chuyển động với tốc độ $20\ \mathrm{m/s}$. Tính động lượng của xe.
\loigiai{
$p=mv=3\cdot20$. Suy ra kết quả $60\ \mathrm{kg\,m/s}$.
}
\end{bt}

% ===== Câu 51 =====
% ID: AIQ_L10C5_FULL_0051
% Mức: VD
\begin{bt}
% ID: AIQ_L10C5_FULL_0051
% Mức: VD
Trình bày cách giải: Xe nhỏ khối lượng $4\ \mathrm{kg}$ chuyển động với tốc độ $21\ \mathrm{m/s}$. Tính động lượng của xe.
\loigiai{
$p=mv=4\cdot21$. Suy ra kết quả $84\ \mathrm{kg\,m/s}$.
}
\end{bt}

% ===== Câu 52 =====
% ID: AIQ_L10C5_FULL_0052
% Mức: VD
\begin{bt}
% ID: AIQ_L10C5_FULL_0052
% Mức: VD
Trình bày cách giải: Xe nhỏ khối lượng $5\ \mathrm{kg}$ chuyển động với tốc độ $22\ \mathrm{m/s}$. Tính động lượng của xe.
\loigiai{
$p=mv=5\cdot22$. Suy ra kết quả $110\ \mathrm{kg\,m/s}$.
}
\end{bt}

% ===== Câu 53 =====
% ID: AIQ_L10C5_FULL_0053
% Mức: VDC
\begin{bt}
% ID: AIQ_L10C5_FULL_0053
% Mức: VDC
Trình bày cách giải: Xe nhỏ khối lượng $2\ \mathrm{kg}$ chuyển động với tốc độ $23\ \mathrm{m/s}$. Tính động lượng của xe.
\loigiai{
$p=mv=2\cdot23$. Suy ra kết quả $46\ \mathrm{kg\,m/s}$.
}
\end{bt}

% ===== Câu 54 =====
% ID: AIQ_L10C5_FULL_0054
% Mức: VDC
\begin{bt}
% ID: AIQ_L10C5_FULL_0054
% Mức: VDC
Trình bày cách giải: Xe nhỏ khối lượng $3\ \mathrm{kg}$ chuyển động với tốc độ $24\ \mathrm{m/s}$. Tính động lượng của xe.
\loigiai{
$p=mv=3\cdot24$. Suy ra kết quả $72\ \mathrm{kg\,m/s}$.
}
\end{bt}

% ===== Câu 55 =====
% ID: AIQ_L10C5_FULL_0055
% Mức: NB
\dangbt{Tổng hợp động lượng của hệ nhiều vật}
\begin{ex}
% ID: AIQ_L10C5_FULL_0055
% Mức: NB
Hai vật có động lượng cùng phương cùng chiều, độ lớn $p_1=6$ và $p_2=2\ \mathrm{kg\,m/s}$. Tính độ lớn tổng động lượng.
\choice
{\True $8\ \mathrm{kg\,m/s}$}
{$4\ \mathrm{kg\,m/s}$}
{$12\ \mathrm{kg\,m/s}$}
{$10\ \mathrm{kg\,m/s}$}
\loigiai{
Cùng chiều nên $p=p_1+p_2=6+2$. Suy ra kết quả $8\ \mathrm{kg\,m/s}$.
}
\end{ex}

% ===== Câu 56 =====
% ID: AIQ_L10C5_FULL_0056
% Mức: NB
\begin{ex}
% ID: AIQ_L10C5_FULL_0056
% Mức: NB
Hai vật có động lượng cùng phương ngược chiều, độ lớn $p_1=7$ và $p_2=3\ \mathrm{kg\,m/s}$. Tính độ lớn tổng động lượng.
\choice
{$10\ \mathrm{kg\,m/s}$}
{\True $4\ \mathrm{kg\,m/s}$}
{$21\ \mathrm{kg\,m/s}$}
{$6\ \mathrm{kg\,m/s}$}
\loigiai{
Ngược chiều nên $p=|p_1-p_2|=|7-3|$. Suy ra kết quả $4\ \mathrm{kg\,m/s}$.
}
\end{ex}

% ===== Câu 57 =====
% ID: AIQ_L10C5_FULL_0057
% Mức: TH
\begin{ex}
% ID: AIQ_L10C5_FULL_0057
% Mức: TH
Hai vật có động lượng cùng phương cùng chiều, độ lớn $p_1=8$ và $p_2=4\ \mathrm{kg\,m/s}$. Tính độ lớn tổng động lượng.
\choice
{$4\ \mathrm{kg\,m/s}$}
{$32\ \mathrm{kg\,m/s}$}
{\True $12\ \mathrm{kg\,m/s}$}
{$14\ \mathrm{kg\,m/s}$}
\loigiai{
Cùng chiều nên $p=p_1+p_2=8+4$. Suy ra kết quả $12\ \mathrm{kg\,m/s}$.
}
\end{ex}

% ===== Câu 58 =====
% ID: AIQ_L10C5_FULL_0058
% Mức: TH
\begin{ex}
% ID: AIQ_L10C5_FULL_0058
% Mức: TH
Hai vật có động lượng cùng phương ngược chiều, độ lớn $p_1=9$ và $p_2=5\ \mathrm{kg\,m/s}$. Tính độ lớn tổng động lượng.
\choice
{$14\ \mathrm{kg\,m/s}$}
{$45\ \mathrm{kg\,m/s}$}
{$6\ \mathrm{kg\,m/s}$}
{\True $4\ \mathrm{kg\,m/s}$}
\loigiai{
Ngược chiều nên $p=|p_1-p_2|=|9-5|$. Suy ra kết quả $4\ \mathrm{kg\,m/s}$.
}
\end{ex}

% ===== Câu 59 =====
% ID: AIQ_L10C5_FULL_0059
% Mức: VD
\begin{ex}
% ID: AIQ_L10C5_FULL_0059
% Mức: VD
Hai vật có động lượng cùng phương cùng chiều, độ lớn $p_1=10$ và $p_2=6\ \mathrm{kg\,m/s}$. Tính độ lớn tổng động lượng.
\choice
{\True $16\ \mathrm{kg\,m/s}$}
{$4\ \mathrm{kg\,m/s}$}
{$60\ \mathrm{kg\,m/s}$}
{$18\ \mathrm{kg\,m/s}$}
\loigiai{
Cùng chiều nên $p=p_1+p_2=10+6$. Suy ra kết quả $16\ \mathrm{kg\,m/s}$.
}
\end{ex}

% ===== Câu 60 =====
% ID: AIQ_L10C5_FULL_0060
% Mức: VD
\begin{ex}
% ID: AIQ_L10C5_FULL_0060
% Mức: VD
Hai vật có động lượng cùng phương ngược chiều, độ lớn $p_1=11$ và $p_2=2\ \mathrm{kg\,m/s}$. Tính độ lớn tổng động lượng.
\choice
{$13\ \mathrm{kg\,m/s}$}
{\True $9\ \mathrm{kg\,m/s}$}
{$22\ \mathrm{kg\,m/s}$}
{$11\ \mathrm{kg\,m/s}$}
\loigiai{
Ngược chiều nên $p=|p_1-p_2|=|11-2|$. Suy ra kết quả $9\ \mathrm{kg\,m/s}$.
}
\end{ex}

% ===== Câu 61 =====
% ID: AIQ_L10C5_FULL_0061
% Mức: VD
\begin{ex}
% ID: AIQ_L10C5_FULL_0061
% Mức: VD
Hai vật có động lượng cùng phương cùng chiều, độ lớn $p_1=12$ và $p_2=3\ \mathrm{kg\,m/s}$. Tính độ lớn tổng động lượng.
\choice
{$9\ \mathrm{kg\,m/s}$}
{$36\ \mathrm{kg\,m/s}$}
{\True $15\ \mathrm{kg\,m/s}$}
{$17\ \mathrm{kg\,m/s}$}
\loigiai{
Cùng chiều nên $p=p_1+p_2=12+3$. Suy ra kết quả $15\ \mathrm{kg\,m/s}$.
}
\end{ex}

% ===== Câu 62 =====
% ID: AIQ_L10C5_FULL_0062
% Mức: VDC
\begin{ex}
% ID: AIQ_L10C5_FULL_0062
% Mức: VDC
Hai vật có động lượng cùng phương ngược chiều, độ lớn $p_1=13$ và $p_2=4\ \mathrm{kg\,m/s}$. Tính độ lớn tổng động lượng.
\choice
{$17\ \mathrm{kg\,m/s}$}
{$52\ \mathrm{kg\,m/s}$}
{$11\ \mathrm{kg\,m/s}$}
{\True $9\ \mathrm{kg\,m/s}$}
\loigiai{
Ngược chiều nên $p=|p_1-p_2|=|13-4|$. Suy ra kết quả $9\ \mathrm{kg\,m/s}$.
}
\end{ex}

% ===== Câu 63 =====
% ID: AIQ_L10C5_FULL_0063
% Mức: VDC
\begin{ex}
% ID: AIQ_L10C5_FULL_0063
% Mức: VDC
Hai vật có động lượng cùng phương cùng chiều, độ lớn $p_1=14$ và $p_2=5\ \mathrm{kg\,m/s}$. Tính độ lớn tổng động lượng.
\choice
{\True $19\ \mathrm{kg\,m/s}$}
{$9\ \mathrm{kg\,m/s}$}
{$70\ \mathrm{kg\,m/s}$}
{$21\ \mathrm{kg\,m/s}$}
\loigiai{
Cùng chiều nên $p=p_1+p_2=14+5$. Suy ra kết quả $19\ \mathrm{kg\,m/s}$.
}
\end{ex}

% ===== Câu 64 =====
% ID: AIQ_L10C5_FULL_0064
% Mức: VDC
\begin{ex}
% ID: AIQ_L10C5_FULL_0064
% Mức: VDC
Hai vật có động lượng cùng phương ngược chiều, độ lớn $p_1=15$ và $p_2=6\ \mathrm{kg\,m/s}$. Tính độ lớn tổng động lượng.
\choice
{$21\ \mathrm{kg\,m/s}$}
{\True $9\ \mathrm{kg\,m/s}$}
{$90\ \mathrm{kg\,m/s}$}
{$11\ \mathrm{kg\,m/s}$}
\loigiai{
Ngược chiều nên $p=|p_1-p_2|=|15-6|$. Suy ra kết quả $9\ \mathrm{kg\,m/s}$.
}
\end{ex}

% ===== Câu 65 =====
% ID: AIQ_L10C5_FULL_0065
% Mức: NB
\begin{ex}
% ID: AIQ_L10C5_FULL_0065
% Mức: NB
Xét các phát biểu thuộc dạng “Tổng hợp động lượng của hệ nhiều vật” ở tình huống số 1.
\choiceTF
{\True Tổng động lượng là tổng vectơ.}
{Hai động lượng ngược chiều luôn cộng độ lớn.}
{\True Hai động lượng vuông góc dùng định lí Pythagore.}
{\True Tổng động lượng có thể bằng 0.}
\loigiai{
\begin{itemchoice}
\item (Đúng) Tổng động lượng là tổng vectơ. (Vì): Phù hợp với định nghĩa hoặc hệ thức vật lí. b) Sai: Không phù hợp với định nghĩa hoặc hệ thức vật lí. c) Đúng: Phù hợp với định nghĩa hoặc hệ thức vật lí. d) Đúng: Phù hợp với định nghĩa hoặc hệ thức vật lí
\item (Sai) Hai động lượng ngược chiều luôn cộng độ lớn.
\item (Đúng) Hai động lượng vuông góc dùng định lí Pythagore.
\item (Đúng) Tổng động lượng có thể bằng 0.
\end{itemchoice}
}
\end{ex}

% ===== Câu 66 =====
% ID: AIQ_L10C5_FULL_0066
% Mức: TH
\begin{ex}
% ID: AIQ_L10C5_FULL_0066
% Mức: TH
Xét các phát biểu thuộc dạng “Tổng hợp động lượng của hệ nhiều vật” ở tình huống số 2.
\choiceTF
{\True Tổng động lượng là tổng vectơ.}
{Hai động lượng ngược chiều luôn cộng độ lớn.}
{\True Hai động lượng vuông góc dùng định lí Pythagore.}
{\True Tổng động lượng có thể bằng 0.}
\loigiai{
\begin{itemchoice}
\item (Đúng) Tổng động lượng là tổng vectơ. (Vì): Phù hợp với định nghĩa hoặc hệ thức vật lí. b) Sai: Không phù hợp với định nghĩa hoặc hệ thức vật lí. c) Đúng: Phù hợp với định nghĩa hoặc hệ thức vật lí. d) Đúng: Phù hợp với định nghĩa hoặc hệ thức vật lí
\item (Sai) Hai động lượng ngược chiều luôn cộng độ lớn.
\item (Đúng) Hai động lượng vuông góc dùng định lí Pythagore.
\item (Đúng) Tổng động lượng có thể bằng 0.
\end{itemchoice}
}
\end{ex}

% ===== Câu 67 =====
% ID: AIQ_L10C5_FULL_0067
% Mức: TH
\begin{ex}
% ID: AIQ_L10C5_FULL_0067
% Mức: TH
Xét các phát biểu thuộc dạng “Tổng hợp động lượng của hệ nhiều vật” ở tình huống số 3.
\choiceTF
{\True Tổng động lượng là tổng vectơ.}
{Hai động lượng ngược chiều luôn cộng độ lớn.}
{\True Hai động lượng vuông góc dùng định lí Pythagore.}
{\True Tổng động lượng có thể bằng 0.}
\loigiai{
\begin{itemchoice}
\item (Đúng) Tổng động lượng là tổng vectơ. (Vì): Phù hợp với định nghĩa hoặc hệ thức vật lí. b) Sai: Không phù hợp với định nghĩa hoặc hệ thức vật lí. c) Đúng: Phù hợp với định nghĩa hoặc hệ thức vật lí. d) Đúng: Phù hợp với định nghĩa hoặc hệ thức vật lí
\item (Sai) Hai động lượng ngược chiều luôn cộng độ lớn.
\item (Đúng) Hai động lượng vuông góc dùng định lí Pythagore.
\item (Đúng) Tổng động lượng có thể bằng 0.
\end{itemchoice}
}
\end{ex}

% ===== Câu 68 =====
% ID: AIQ_L10C5_FULL_0068
% Mức: VD
\begin{ex}
% ID: AIQ_L10C5_FULL_0068
% Mức: VD
Xét các phát biểu thuộc dạng “Tổng hợp động lượng của hệ nhiều vật” ở tình huống số 4.
\choiceTF
{\True Tổng động lượng là tổng vectơ.}
{Hai động lượng ngược chiều luôn cộng độ lớn.}
{\True Hai động lượng vuông góc dùng định lí Pythagore.}
{\True Tổng động lượng có thể bằng 0.}
\loigiai{
\begin{itemchoice}
\item (Đúng) Tổng động lượng là tổng vectơ. (Vì): Phù hợp với định nghĩa hoặc hệ thức vật lí. b) Sai: Không phù hợp với định nghĩa hoặc hệ thức vật lí. c) Đúng: Phù hợp với định nghĩa hoặc hệ thức vật lí. d) Đúng: Phù hợp với định nghĩa hoặc hệ thức vật lí
\item (Sai) Hai động lượng ngược chiều luôn cộng độ lớn.
\item (Đúng) Hai động lượng vuông góc dùng định lí Pythagore.
\item (Đúng) Tổng động lượng có thể bằng 0.
\end{itemchoice}
}
\end{ex}

% ===== Câu 69 =====
% ID: AIQ_L10C5_FULL_0069
% Mức: VDC
\begin{ex}
% ID: AIQ_L10C5_FULL_0069
% Mức: VDC
Xét các phát biểu thuộc dạng “Tổng hợp động lượng của hệ nhiều vật” ở tình huống số 5.
\choiceTF
{\True Tổng động lượng là tổng vectơ.}
{Hai động lượng ngược chiều luôn cộng độ lớn.}
{\True Hai động lượng vuông góc dùng định lí Pythagore.}
{\True Tổng động lượng có thể bằng 0.}
\loigiai{
\begin{itemchoice}
\item (Đúng) Tổng động lượng là tổng vectơ. (Vì): Phù hợp với định nghĩa hoặc hệ thức vật lí. b) Sai: Không phù hợp với định nghĩa hoặc hệ thức vật lí. c) Đúng: Phù hợp với định nghĩa hoặc hệ thức vật lí. d) Đúng: Phù hợp với định nghĩa hoặc hệ thức vật lí
\item (Sai) Hai động lượng ngược chiều luôn cộng độ lớn.
\item (Đúng) Hai động lượng vuông góc dùng định lí Pythagore.
\item (Đúng) Tổng động lượng có thể bằng 0.
\end{itemchoice}
}
\end{ex}

% ===== Câu 70 =====
% ID: AIQ_L10C5_FULL_0070
% Mức: TH
\begin{ex}
% ID: AIQ_L10C5_FULL_0070
% Mức: TH
Hai vật có động lượng cùng phương cùng chiều, độ lớn $p_1=16$ và $p_2=2\ \mathrm{kg\,m/s}$. Tính độ lớn tổng động lượng.
\shortans{18}
\loigiai{
Cùng chiều nên $p=p_1+p_2=16+2$. Suy ra kết quả $18\ \mathrm{kg\,m/s}$. Đáp án trả lời ngắn không ghi đơn vị.
}
\end{ex}

% ===== Câu 71 =====
% ID: AIQ_L10C5_FULL_0071
% Mức: TH
\begin{ex}
% ID: AIQ_L10C5_FULL_0071
% Mức: TH
Hai vật có động lượng cùng phương ngược chiều, độ lớn $p_1=17$ và $p_2=3\ \mathrm{kg\,m/s}$. Tính độ lớn tổng động lượng.
\shortans{14}
\loigiai{
Ngược chiều nên $p=|p_1-p_2|=|17-3|$. Suy ra kết quả $14\ \mathrm{kg\,m/s}$. Đáp án trả lời ngắn không ghi đơn vị.
}
\end{ex}

% ===== Câu 72 =====
% ID: AIQ_L10C5_FULL_0072
% Mức: VD
\begin{ex}
% ID: AIQ_L10C5_FULL_0072
% Mức: VD
Hai vật có động lượng cùng phương cùng chiều, độ lớn $p_1=18$ và $p_2=4\ \mathrm{kg\,m/s}$. Tính độ lớn tổng động lượng.
\shortans{22}
\loigiai{
Cùng chiều nên $p=p_1+p_2=18+4$. Suy ra kết quả $22\ \mathrm{kg\,m/s}$. Đáp án trả lời ngắn không ghi đơn vị.
}
\end{ex}

% ===== Câu 73 =====
% ID: AIQ_L10C5_FULL_0073
% Mức: VD
\begin{ex}
% ID: AIQ_L10C5_FULL_0073
% Mức: VD
Hai vật có động lượng cùng phương ngược chiều, độ lớn $p_1=19$ và $p_2=5\ \mathrm{kg\,m/s}$. Tính độ lớn tổng động lượng.
\shortans{14}
\loigiai{
Ngược chiều nên $p=|p_1-p_2|=|19-5|$. Suy ra kết quả $14\ \mathrm{kg\,m/s}$. Đáp án trả lời ngắn không ghi đơn vị.
}
\end{ex}

% ===== Câu 74 =====
% ID: AIQ_L10C5_FULL_0074
% Mức: VDC
\begin{ex}
% ID: AIQ_L10C5_FULL_0074
% Mức: VDC
Hai vật có động lượng cùng phương cùng chiều, độ lớn $p_1=20$ và $p_2=6\ \mathrm{kg\,m/s}$. Tính độ lớn tổng động lượng.
\shortans{26}
\loigiai{
Cùng chiều nên $p=p_1+p_2=20+6$. Suy ra kết quả $26\ \mathrm{kg\,m/s}$. Đáp án trả lời ngắn không ghi đơn vị.
}
\end{ex}

% ===== Câu 75 =====
% ID: AIQ_L10C5_FULL_0075
% Mức: VDC
\begin{ex}
% ID: AIQ_L10C5_FULL_0075
% Mức: VDC
Hai vật có động lượng cùng phương ngược chiều, độ lớn $p_1=21$ và $p_2=2\ \mathrm{kg\,m/s}$. Tính độ lớn tổng động lượng.
\shortans{19}
\loigiai{
Ngược chiều nên $p=|p_1-p_2|=|21-2|$. Suy ra kết quả $19\ \mathrm{kg\,m/s}$. Đáp án trả lời ngắn không ghi đơn vị.
}
\end{ex}

% ===== Câu 76 =====
% ID: AIQ_L10C5_FULL_0076
% Mức: TH
\begin{bt}
% ID: AIQ_L10C5_FULL_0076
% Mức: TH
Trình bày cách giải: Hai vật có động lượng cùng phương cùng chiều, độ lớn $p_1=22$ và $p_2=3\ \mathrm{kg\,m/s}$. Tính độ lớn tổng động lượng.
\loigiai{
Cùng chiều nên $p=p_1+p_2=22+3$. Suy ra kết quả $25\ \mathrm{kg\,m/s}$.
}
\end{bt}

% ===== Câu 77 =====
% ID: AIQ_L10C5_FULL_0077
% Mức: TH
\begin{bt}
% ID: AIQ_L10C5_FULL_0077
% Mức: TH
Trình bày cách giải: Hai vật có động lượng cùng phương ngược chiều, độ lớn $p_1=23$ và $p_2=4\ \mathrm{kg\,m/s}$. Tính độ lớn tổng động lượng.
\loigiai{
Ngược chiều nên $p=|p_1-p_2|=|23-4|$. Suy ra kết quả $19\ \mathrm{kg\,m/s}$.
}
\end{bt}

% ===== Câu 78 =====
% ID: AIQ_L10C5_FULL_0078
% Mức: VD
\begin{bt}
% ID: AIQ_L10C5_FULL_0078
% Mức: VD
Trình bày cách giải: Hai vật có động lượng cùng phương cùng chiều, độ lớn $p_1=24$ và $p_2=5\ \mathrm{kg\,m/s}$. Tính độ lớn tổng động lượng.
\loigiai{
Cùng chiều nên $p=p_1+p_2=24+5$. Suy ra kết quả $29\ \mathrm{kg\,m/s}$.
}
\end{bt}

% ===== Câu 79 =====
% ID: AIQ_L10C5_FULL_0079
% Mức: VD
\begin{bt}
% ID: AIQ_L10C5_FULL_0079
% Mức: VD
Trình bày cách giải: Hai vật có động lượng cùng phương ngược chiều, độ lớn $p_1=25$ và $p_2=6\ \mathrm{kg\,m/s}$. Tính độ lớn tổng động lượng.
\loigiai{
Ngược chiều nên $p=|p_1-p_2|=|25-6|$. Suy ra kết quả $19\ \mathrm{kg\,m/s}$.
}
\end{bt}

% ===== Câu 80 =====
% ID: AIQ_L10C5_FULL_0080
% Mức: VDC
\begin{bt}
% ID: AIQ_L10C5_FULL_0080
% Mức: VDC
Trình bày cách giải: Hai vật có động lượng cùng phương cùng chiều, độ lớn $p_1=26$ và $p_2=2\ \mathrm{kg\,m/s}$. Tính độ lớn tổng động lượng.
\loigiai{
Cùng chiều nên $p=p_1+p_2=26+2$. Suy ra kết quả $28\ \mathrm{kg\,m/s}$.
}
\end{bt}

% ===== Câu 81 =====
% ID: AIQ_L10C5_FULL_0081
% Mức: VDC
\begin{bt}
% ID: AIQ_L10C5_FULL_0081
% Mức: VDC
Trình bày cách giải: Hai vật có động lượng cùng phương ngược chiều, độ lớn $p_1=27$ và $p_2=3\ \mathrm{kg\,m/s}$. Tính độ lớn tổng động lượng.
\loigiai{
Ngược chiều nên $p=|p_1-p_2|=|27-3|$. Suy ra kết quả $24\ \mathrm{kg\,m/s}$.
}
\end{bt}

% ===== Câu 82 =====
% ID: AIQ_L10C5_FULL_0082
% Mức: NB
\dangbt{Liên hệ độ biến thiên động lượng với xung lượng của lực}
\begin{ex}
% ID: AIQ_L10C5_FULL_0082
% Mức: NB
Một lực không đổi $10\ \mathrm{N}$ tác dụng lên vật trong $0,2\ \mathrm{s}$. Tính độ lớn xung lượng của lực.
\choice
{\True $2\ \mathrm{N\,s}$}
{$10,2\ \mathrm{N\,s}$}
{$50\ \mathrm{N\,s}$}
{$20\ \mathrm{N\,s}$}
\loigiai{
$J=F\Delta t=10\cdot0,2$. Suy ra kết quả $2\ \mathrm{N\,s}$.
}
\end{ex}

% ===== Câu 83 =====
% ID: AIQ_L10C5_FULL_0083
% Mức: NB
\begin{ex}
% ID: AIQ_L10C5_FULL_0083
% Mức: NB
Một lực không đổi $12\ \mathrm{N}$ tác dụng lên vật trong $0,3\ \mathrm{s}$. Tính độ lớn xung lượng của lực.
\choice
{$12,3\ \mathrm{N\,s}$}
{\True $3,6\ \mathrm{N\,s}$}
{$40\ \mathrm{N\,s}$}
{$36\ \mathrm{N\,s}$}
\loigiai{
$J=F\Delta t=12\cdot0,3$. Suy ra kết quả $3,6\ \mathrm{N\,s}$.
}
\end{ex}

% ===== Câu 84 =====
% ID: AIQ_L10C5_FULL_0084
% Mức: TH
\begin{ex}
% ID: AIQ_L10C5_FULL_0084
% Mức: TH
Một lực không đổi $14\ \mathrm{N}$ tác dụng lên vật trong $0,4\ \mathrm{s}$. Tính độ lớn xung lượng của lực.
\choice
{$14,4\ \mathrm{N\,s}$}
{$35\ \mathrm{N\,s}$}
{\True $5,6\ \mathrm{N\,s}$}
{$56\ \mathrm{N\,s}$}
\loigiai{
$J=F\Delta t=14\cdot0,4$. Suy ra kết quả $5,6\ \mathrm{N\,s}$.
}
\end{ex}

% ===== Câu 85 =====
% ID: AIQ_L10C5_FULL_0085
% Mức: TH
\begin{ex}
% ID: AIQ_L10C5_FULL_0085
% Mức: TH
Một lực không đổi $16\ \mathrm{N}$ tác dụng lên vật trong $0,5\ \mathrm{s}$. Tính độ lớn xung lượng của lực.
\choice
{$16,5\ \mathrm{N\,s}$}
{$32\ \mathrm{N\,s}$}
{$80\ \mathrm{N\,s}$}
{\True $8\ \mathrm{N\,s}$}
\loigiai{
$J=F\Delta t=16\cdot0,5$. Suy ra kết quả $8\ \mathrm{N\,s}$.
}
\end{ex}

% ===== Câu 86 =====
% ID: AIQ_L10C5_FULL_0086
% Mức: VD
\begin{ex}
% ID: AIQ_L10C5_FULL_0086
% Mức: VD
Một lực không đổi $18\ \mathrm{N}$ tác dụng lên vật trong $0,6\ \mathrm{s}$. Tính độ lớn xung lượng của lực.
\choice
{\True $10,8\ \mathrm{N\,s}$}
{$18,6\ \mathrm{N\,s}$}
{$30\ \mathrm{N\,s}$}
{$108\ \mathrm{N\,s}$}
\loigiai{
$J=F\Delta t=18\cdot0,6$. Suy ra kết quả $10,8\ \mathrm{N\,s}$.
}
\end{ex}

% ===== Câu 87 =====
% ID: AIQ_L10C5_FULL_0087
% Mức: VD
\begin{ex}
% ID: AIQ_L10C5_FULL_0087
% Mức: VD
Một lực không đổi $20\ \mathrm{N}$ tác dụng lên vật trong $0,2\ \mathrm{s}$. Tính độ lớn xung lượng của lực.
\choice
{$20,2\ \mathrm{N\,s}$}
{\True $4\ \mathrm{N\,s}$}
{$100\ \mathrm{N\,s}$}
{$40\ \mathrm{N\,s}$}
\loigiai{
$J=F\Delta t=20\cdot0,2$. Suy ra kết quả $4\ \mathrm{N\,s}$.
}
\end{ex}

% ===== Câu 88 =====
% ID: AIQ_L10C5_FULL_0088
% Mức: VD
\begin{ex}
% ID: AIQ_L10C5_FULL_0088
% Mức: VD
Một lực không đổi $22\ \mathrm{N}$ tác dụng lên vật trong $0,3\ \mathrm{s}$. Tính độ lớn xung lượng của lực.
\choice
{$22,3\ \mathrm{N\,s}$}
{$73,33\ \mathrm{N\,s}$}
{\True $6,6\ \mathrm{N\,s}$}
{$66\ \mathrm{N\,s}$}
\loigiai{
$J=F\Delta t=22\cdot0,3$. Suy ra kết quả $6,6\ \mathrm{N\,s}$.
}
\end{ex}

% ===== Câu 89 =====
% ID: AIQ_L10C5_FULL_0089
% Mức: VDC
\begin{ex}
% ID: AIQ_L10C5_FULL_0089
% Mức: VDC
Một lực không đổi $24\ \mathrm{N}$ tác dụng lên vật trong $0,4\ \mathrm{s}$. Tính độ lớn xung lượng của lực.
\choice
{$24,4\ \mathrm{N\,s}$}
{$60\ \mathrm{N\,s}$}
{$96\ \mathrm{N\,s}$}
{\True $9,6\ \mathrm{N\,s}$}
\loigiai{
$J=F\Delta t=24\cdot0,4$. Suy ra kết quả $9,6\ \mathrm{N\,s}$.
}
\end{ex}

% ===== Câu 90 =====
% ID: AIQ_L10C5_FULL_0090
% Mức: VDC
\begin{ex}
% ID: AIQ_L10C5_FULL_0090
% Mức: VDC
Một lực không đổi $26\ \mathrm{N}$ tác dụng lên vật trong $0,5\ \mathrm{s}$. Tính độ lớn xung lượng của lực.
\choice
{\True $13\ \mathrm{N\,s}$}
{$26,5\ \mathrm{N\,s}$}
{$52\ \mathrm{N\,s}$}
{$130\ \mathrm{N\,s}$}
\loigiai{
$J=F\Delta t=26\cdot0,5$. Suy ra kết quả $13\ \mathrm{N\,s}$.
}
\end{ex}

% ===== Câu 91 =====
% ID: AIQ_L10C5_FULL_0091
% Mức: VDC
\begin{ex}
% ID: AIQ_L10C5_FULL_0091
% Mức: VDC
Một lực không đổi $28\ \mathrm{N}$ tác dụng lên vật trong $0,6\ \mathrm{s}$. Tính độ lớn xung lượng của lực.
\choice
{$28,6\ \mathrm{N\,s}$}
{\True $16,8\ \mathrm{N\,s}$}
{$46,67\ \mathrm{N\,s}$}
{$168\ \mathrm{N\,s}$}
\loigiai{
$J=F\Delta t=28\cdot0,6$. Suy ra kết quả $16,8\ \mathrm{N\,s}$.
}
\end{ex}

% ===== Câu 92 =====
% ID: AIQ_L10C5_FULL_0092
% Mức: NB
\begin{ex}
% ID: AIQ_L10C5_FULL_0092
% Mức: NB
Xét các phát biểu thuộc dạng “Liên hệ độ biến thiên động lượng với xung lượng của lực” ở tình huống số 1.
\choiceTF
{\True Xung lượng bằng $F\Delta t$ khi lực không đổi.}
{\True Xung lượng bằng độ biến thiên động lượng.}
{Đơn vị xung lượng là N/s.}
{\True Tăng thời gian hãm có thể giảm lực trung bình.}
\loigiai{
\begin{itemchoice}
\item (Đúng) Xung lượng bằng $F\Delta t$ khi lực không đổi. (Vì): Phù hợp với định nghĩa hoặc hệ thức vật lí. b) Đúng: Phù hợp với định nghĩa hoặc hệ thức vật lí. c) Sai: Không phù hợp với định nghĩa hoặc hệ thức vật lí. d) Đúng: Phù hợp với định nghĩa hoặc hệ thức vật lí
\item (Đúng) Xung lượng bằng độ biến thiên động lượng.
\item (Sai) Đơn vị xung lượng là N/s.
\item (Đúng) Tăng thời gian hãm có thể giảm lực trung bình.
\end{itemchoice}
}
\end{ex}

% ===== Câu 93 =====
% ID: AIQ_L10C5_FULL_0093
% Mức: TH
\begin{ex}
% ID: AIQ_L10C5_FULL_0093
% Mức: TH
Xét các phát biểu thuộc dạng “Liên hệ độ biến thiên động lượng với xung lượng của lực” ở tình huống số 2.
\choiceTF
{\True Xung lượng bằng $F\Delta t$ khi lực không đổi.}
{\True Xung lượng bằng độ biến thiên động lượng.}
{Đơn vị xung lượng là N/s.}
{\True Tăng thời gian hãm có thể giảm lực trung bình.}
\loigiai{
\begin{itemchoice}
\item (Đúng) Xung lượng bằng $F\Delta t$ khi lực không đổi. (Vì): Phù hợp với định nghĩa hoặc hệ thức vật lí. b) Đúng: Phù hợp với định nghĩa hoặc hệ thức vật lí. c) Sai: Không phù hợp với định nghĩa hoặc hệ thức vật lí. d) Đúng: Phù hợp với định nghĩa hoặc hệ thức vật lí
\item (Đúng) Xung lượng bằng độ biến thiên động lượng.
\item (Sai) Đơn vị xung lượng là N/s.
\item (Đúng) Tăng thời gian hãm có thể giảm lực trung bình.
\end{itemchoice}
}
\end{ex}

% ===== Câu 94 =====
% ID: AIQ_L10C5_FULL_0094
% Mức: TH
\begin{ex}
% ID: AIQ_L10C5_FULL_0094
% Mức: TH
Xét các phát biểu thuộc dạng “Liên hệ độ biến thiên động lượng với xung lượng của lực” ở tình huống số 3.
\choiceTF
{\True Xung lượng bằng $F\Delta t$ khi lực không đổi.}
{\True Xung lượng bằng độ biến thiên động lượng.}
{Đơn vị xung lượng là N/s.}
{\True Tăng thời gian hãm có thể giảm lực trung bình.}
\loigiai{
\begin{itemchoice}
\item (Đúng) Xung lượng bằng $F\Delta t$ khi lực không đổi. (Vì): Phù hợp với định nghĩa hoặc hệ thức vật lí. b) Đúng: Phù hợp với định nghĩa hoặc hệ thức vật lí. c) Sai: Không phù hợp với định nghĩa hoặc hệ thức vật lí. d) Đúng: Phù hợp với định nghĩa hoặc hệ thức vật lí
\item (Đúng) Xung lượng bằng độ biến thiên động lượng.
\item (Sai) Đơn vị xung lượng là N/s.
\item (Đúng) Tăng thời gian hãm có thể giảm lực trung bình.
\end{itemchoice}
}
\end{ex}

% ===== Câu 95 =====
% ID: AIQ_L10C5_FULL_0095
% Mức: VD
\begin{ex}
% ID: AIQ_L10C5_FULL_0095
% Mức: VD
Xét các phát biểu thuộc dạng “Liên hệ độ biến thiên động lượng với xung lượng của lực” ở tình huống số 4.
\choiceTF
{\True Xung lượng bằng $F\Delta t$ khi lực không đổi.}
{\True Xung lượng bằng độ biến thiên động lượng.}
{Đơn vị xung lượng là N/s.}
{\True Tăng thời gian hãm có thể giảm lực trung bình.}
\loigiai{
\begin{itemchoice}
\item (Đúng) Xung lượng bằng $F\Delta t$ khi lực không đổi. (Vì): Phù hợp với định nghĩa hoặc hệ thức vật lí. b) Đúng: Phù hợp với định nghĩa hoặc hệ thức vật lí. c) Sai: Không phù hợp với định nghĩa hoặc hệ thức vật lí. d) Đúng: Phù hợp với định nghĩa hoặc hệ thức vật lí
\item (Đúng) Xung lượng bằng độ biến thiên động lượng.
\item (Sai) Đơn vị xung lượng là N/s.
\item (Đúng) Tăng thời gian hãm có thể giảm lực trung bình.
\end{itemchoice}
}
\end{ex}

% ===== Câu 96 =====
% ID: AIQ_L10C5_FULL_0096
% Mức: VDC
\begin{ex}
% ID: AIQ_L10C5_FULL_0096
% Mức: VDC
Xét các phát biểu thuộc dạng “Liên hệ độ biến thiên động lượng với xung lượng của lực” ở tình huống số 5.
\choiceTF
{\True Xung lượng bằng $F\Delta t$ khi lực không đổi.}
{\True Xung lượng bằng độ biến thiên động lượng.}
{Đơn vị xung lượng là N/s.}
{\True Tăng thời gian hãm có thể giảm lực trung bình.}
\loigiai{
\begin{itemchoice}
\item (Đúng) Xung lượng bằng $F\Delta t$ khi lực không đổi. (Vì): Phù hợp với định nghĩa hoặc hệ thức vật lí. b) Đúng: Phù hợp với định nghĩa hoặc hệ thức vật lí. c) Sai: Không phù hợp với định nghĩa hoặc hệ thức vật lí. d) Đúng: Phù hợp với định nghĩa hoặc hệ thức vật lí
\item (Đúng) Xung lượng bằng độ biến thiên động lượng.
\item (Sai) Đơn vị xung lượng là N/s.
\item (Đúng) Tăng thời gian hãm có thể giảm lực trung bình.
\end{itemchoice}
}
\end{ex}

% ===== Câu 97 =====
% ID: AIQ_L10C5_FULL_0097
% Mức: TH
\begin{ex}
% ID: AIQ_L10C5_FULL_0097
% Mức: TH
Một lực không đổi $30\ \mathrm{N}$ tác dụng lên vật trong $0,2\ \mathrm{s}$. Tính độ lớn xung lượng của lực.
\shortans{6}
\loigiai{
$J=F\Delta t=30\cdot0,2$. Suy ra kết quả $6\ \mathrm{N\,s}$. Đáp án trả lời ngắn không ghi đơn vị.
}
\end{ex}

% ===== Câu 98 =====
% ID: AIQ_L10C5_FULL_0098
% Mức: TH
\begin{ex}
% ID: AIQ_L10C5_FULL_0098
% Mức: TH
Một lực không đổi $32\ \mathrm{N}$ tác dụng lên vật trong $0,3\ \mathrm{s}$. Tính độ lớn xung lượng của lực.
\shortans{9,6}
\loigiai{
$J=F\Delta t=32\cdot0,3$. Suy ra kết quả $9,6\ \mathrm{N\,s}$. Đáp án trả lời ngắn không ghi đơn vị.
}
\end{ex}

% ===== Câu 99 =====
% ID: AIQ_L10C5_FULL_0099
% Mức: VD
\begin{ex}
% ID: AIQ_L10C5_FULL_0099
% Mức: VD
Một lực không đổi $34\ \mathrm{N}$ tác dụng lên vật trong $0,4\ \mathrm{s}$. Tính độ lớn xung lượng của lực.
\shortans{13,6}
\loigiai{
$J=F\Delta t=34\cdot0,4$. Suy ra kết quả $13,6\ \mathrm{N\,s}$. Đáp án trả lời ngắn không ghi đơn vị.
}
\end{ex}

% ===== Câu 100 =====
% ID: AIQ_L10C5_FULL_0100
% Mức: VD
\begin{ex}
% ID: AIQ_L10C5_FULL_0100
% Mức: VD
Một lực không đổi $36\ \mathrm{N}$ tác dụng lên vật trong $0,5\ \mathrm{s}$. Tính độ lớn xung lượng của lực.
\shortans{18}
\loigiai{
$J=F\Delta t=36\cdot0,5$. Suy ra kết quả $18\ \mathrm{N\,s}$. Đáp án trả lời ngắn không ghi đơn vị.
}
\end{ex}

% ===== Câu 101 =====
% ID: AIQ_L10C5_FULL_0101
% Mức: VDC
\begin{ex}
% ID: AIQ_L10C5_FULL_0101
% Mức: VDC
Một lực không đổi $38\ \mathrm{N}$ tác dụng lên vật trong $0,6\ \mathrm{s}$. Tính độ lớn xung lượng của lực.
\shortans{22,8}
\loigiai{
$J=F\Delta t=38\cdot0,6$. Suy ra kết quả $22,8\ \mathrm{N\,s}$. Đáp án trả lời ngắn không ghi đơn vị.
}
\end{ex}

% ===== Câu 102 =====
% ID: AIQ_L10C5_FULL_0102
% Mức: VDC
\begin{ex}
% ID: AIQ_L10C5_FULL_0102
% Mức: VDC
Một lực không đổi $40\ \mathrm{N}$ tác dụng lên vật trong $0,2\ \mathrm{s}$. Tính độ lớn xung lượng của lực.
\shortans{8}
\loigiai{
$J=F\Delta t=40\cdot0,2$. Suy ra kết quả $8\ \mathrm{N\,s}$. Đáp án trả lời ngắn không ghi đơn vị.
}
\end{ex}

% ===== Câu 103 =====
% ID: AIQ_L10C5_FULL_0103
% Mức: TH
\begin{bt}
% ID: AIQ_L10C5_FULL_0103
% Mức: TH
Trình bày cách giải: Một lực không đổi $42\ \mathrm{N}$ tác dụng lên vật trong $0,3\ \mathrm{s}$. Tính độ lớn xung lượng của lực.
\loigiai{
$J=F\Delta t=42\cdot0,3$. Suy ra kết quả $12,6\ \mathrm{N\,s}$.
}
\end{bt}

% ===== Câu 104 =====
% ID: AIQ_L10C5_FULL_0104
% Mức: TH
\begin{bt}
% ID: AIQ_L10C5_FULL_0104
% Mức: TH
Trình bày cách giải: Một lực không đổi $44\ \mathrm{N}$ tác dụng lên vật trong $0,4\ \mathrm{s}$. Tính độ lớn xung lượng của lực.
\loigiai{
$J=F\Delta t=44\cdot0,4$. Suy ra kết quả $17,6\ \mathrm{N\,s}$.
}
\end{bt}

% ===== Câu 105 =====
% ID: AIQ_L10C5_FULL_0105
% Mức: VD
\begin{bt}
% ID: AIQ_L10C5_FULL_0105
% Mức: VD
Trình bày cách giải: Một lực không đổi $46\ \mathrm{N}$ tác dụng lên vật trong $0,5\ \mathrm{s}$. Tính độ lớn xung lượng của lực.
\loigiai{
$J=F\Delta t=46\cdot0,5$. Suy ra kết quả $23\ \mathrm{N\,s}$.
}
\end{bt}

% ===== Câu 106 =====
% ID: AIQ_L10C5_FULL_0106
% Mức: VD
\begin{bt}
% ID: AIQ_L10C5_FULL_0106
% Mức: VD
Trình bày cách giải: Một lực không đổi $48\ \mathrm{N}$ tác dụng lên vật trong $0,6\ \mathrm{s}$. Tính độ lớn xung lượng của lực.
\loigiai{
$J=F\Delta t=48\cdot0,6$. Suy ra kết quả $28,8\ \mathrm{N\,s}$.
}
\end{bt}

% ===== Câu 107 =====
% ID: AIQ_L10C5_FULL_0107
% Mức: VDC
\begin{bt}
% ID: AIQ_L10C5_FULL_0107
% Mức: VDC
Trình bày cách giải: Một lực không đổi $50\ \mathrm{N}$ tác dụng lên vật trong $0,2\ \mathrm{s}$. Tính độ lớn xung lượng của lực.
\loigiai{
$J=F\Delta t=50\cdot0,2$. Suy ra kết quả $10\ \mathrm{N\,s}$.
}
\end{bt}

% ===== Câu 108 =====
% ID: AIQ_L10C5_FULL_0108
% Mức: VDC
\begin{bt}
% ID: AIQ_L10C5_FULL_0108
% Mức: VDC
Trình bày cách giải: Một lực không đổi $52\ \mathrm{N}$ tác dụng lên vật trong $0,3\ \mathrm{s}$. Tính độ lớn xung lượng của lực.
\loigiai{
$J=F\Delta t=52\cdot0,3$. Suy ra kết quả $15,6\ \mathrm{N\,s}$.
}
\end{bt}

% ===== Câu 109 =====
% ID: AIQ_L10C5_FULL_0109
% Mức: NB
\dangbt{Khai thác đồ thị lực – thời gian}
\begin{ex}
% ID: AIQ_L10C5_FULL_0109
% Mức: NB
Đồ thị $F-t$ là đoạn nằm ngang ở $F=8\ \mathrm{N}$ từ $t=0$ đến $t=1\ \mathrm{s}$. Tính xung lượng của lực.
\choice
{\True $8\ \mathrm{N\,s}$}
{$9\ \mathrm{N\,s}$}
{Không đủ dữ kiện để có giá trị này}
{$4\ \mathrm{N\,s}$}
\loigiai{
Xung lượng bằng diện tích dưới đồ thị: $J=F\Delta t=8\cdot1$. Suy ra kết quả $8\ \mathrm{N\,s}$.
}
\end{ex}

% ===== Câu 110 =====
% ID: AIQ_L10C5_FULL_0110
% Mức: NB
\begin{ex}
% ID: AIQ_L10C5_FULL_0110
% Mức: NB
Đồ thị $F-t$ là đoạn nằm ngang ở $F=10\ \mathrm{N}$ từ $t=0$ đến $t=1,5\ \mathrm{s}$. Tính xung lượng của lực.
\choice
{$11,5\ \mathrm{N\,s}$}
{\True $15\ \mathrm{N\,s}$}
{$6,67\ \mathrm{N\,s}$}
{$7,5\ \mathrm{N\,s}$}
\loigiai{
Xung lượng bằng diện tích dưới đồ thị: $J=F\Delta t=10\cdot1,5$. Suy ra kết quả $15\ \mathrm{N\,s}$.
}
\end{ex}

% ===== Câu 111 =====
% ID: AIQ_L10C5_FULL_0111
% Mức: TH
\begin{ex}
% ID: AIQ_L10C5_FULL_0111
% Mức: TH
Đồ thị $F-t$ là đoạn nằm ngang ở $F=12\ \mathrm{N}$ từ $t=0$ đến $t=2\ \mathrm{s}$. Tính xung lượng của lực.
\choice
{$14\ \mathrm{N\,s}$}
{$6\ \mathrm{N\,s}$}
{\True $24\ \mathrm{N\,s}$}
{$12\ \mathrm{N\,s}$}
\loigiai{
Xung lượng bằng diện tích dưới đồ thị: $J=F\Delta t=12\cdot2$. Suy ra kết quả $24\ \mathrm{N\,s}$.
}
\end{ex}

% ===== Câu 112 =====
% ID: AIQ_L10C5_FULL_0112
% Mức: TH
\begin{ex}
% ID: AIQ_L10C5_FULL_0112
% Mức: TH
Đồ thị $F-t$ là đoạn nằm ngang ở $F=14\ \mathrm{N}$ từ $t=0$ đến $t=2,5\ \mathrm{s}$. Tính xung lượng của lực.
\choice
{$16,5\ \mathrm{N\,s}$}
{$5,6\ \mathrm{N\,s}$}
{$17,5\ \mathrm{N\,s}$}
{\True $35\ \mathrm{N\,s}$}
\loigiai{
Xung lượng bằng diện tích dưới đồ thị: $J=F\Delta t=14\cdot2,5$. Suy ra kết quả $35\ \mathrm{N\,s}$.
}
\end{ex}

% ===== Câu 113 =====
% ID: AIQ_L10C5_FULL_0113
% Mức: VD
\begin{ex}
% ID: AIQ_L10C5_FULL_0113
% Mức: VD
Đồ thị $F-t$ là đoạn nằm ngang ở $F=16\ \mathrm{N}$ từ $t=0$ đến $t=1\ \mathrm{s}$. Tính xung lượng của lực.
\choice
{\True $16\ \mathrm{N\,s}$}
{$17\ \mathrm{N\,s}$}
{Không đủ dữ kiện để có giá trị này}
{$8\ \mathrm{N\,s}$}
\loigiai{
Xung lượng bằng diện tích dưới đồ thị: $J=F\Delta t=16\cdot1$. Suy ra kết quả $16\ \mathrm{N\,s}$.
}
\end{ex}

% ===== Câu 114 =====
% ID: AIQ_L10C5_FULL_0114
% Mức: VD
\begin{ex}
% ID: AIQ_L10C5_FULL_0114
% Mức: VD
Đồ thị $F-t$ là đoạn nằm ngang ở $F=18\ \mathrm{N}$ từ $t=0$ đến $t=1,5\ \mathrm{s}$. Tính xung lượng của lực.
\choice
{$19,5\ \mathrm{N\,s}$}
{\True $27\ \mathrm{N\,s}$}
{$12\ \mathrm{N\,s}$}
{$13,5\ \mathrm{N\,s}$}
\loigiai{
Xung lượng bằng diện tích dưới đồ thị: $J=F\Delta t=18\cdot1,5$. Suy ra kết quả $27\ \mathrm{N\,s}$.
}
\end{ex}

% ===== Câu 115 =====
% ID: AIQ_L10C5_FULL_0115
% Mức: VD
\begin{ex}
% ID: AIQ_L10C5_FULL_0115
% Mức: VD
Đồ thị $F-t$ là đoạn nằm ngang ở $F=20\ \mathrm{N}$ từ $t=0$ đến $t=2\ \mathrm{s}$. Tính xung lượng của lực.
\choice
{$22\ \mathrm{N\,s}$}
{$10\ \mathrm{N\,s}$}
{\True $40\ \mathrm{N\,s}$}
{$20\ \mathrm{N\,s}$}
\loigiai{
Xung lượng bằng diện tích dưới đồ thị: $J=F\Delta t=20\cdot2$. Suy ra kết quả $40\ \mathrm{N\,s}$.
}
\end{ex}

% ===== Câu 116 =====
% ID: AIQ_L10C5_FULL_0116
% Mức: VDC
\begin{ex}
% ID: AIQ_L10C5_FULL_0116
% Mức: VDC
Đồ thị $F-t$ là đoạn nằm ngang ở $F=22\ \mathrm{N}$ từ $t=0$ đến $t=2,5\ \mathrm{s}$. Tính xung lượng của lực.
\choice
{$24,5\ \mathrm{N\,s}$}
{$8,8\ \mathrm{N\,s}$}
{$27,5\ \mathrm{N\,s}$}
{\True $55\ \mathrm{N\,s}$}
\loigiai{
Xung lượng bằng diện tích dưới đồ thị: $J=F\Delta t=22\cdot2,5$. Suy ra kết quả $55\ \mathrm{N\,s}$.
}
\end{ex}

% ===== Câu 117 =====
% ID: AIQ_L10C5_FULL_0117
% Mức: VDC
\begin{ex}
% ID: AIQ_L10C5_FULL_0117
% Mức: VDC
Đồ thị $F-t$ là đoạn nằm ngang ở $F=24\ \mathrm{N}$ từ $t=0$ đến $t=1\ \mathrm{s}$. Tính xung lượng của lực.
\choice
{\True $24\ \mathrm{N\,s}$}
{$25\ \mathrm{N\,s}$}
{Không đủ dữ kiện để có giá trị này}
{$12\ \mathrm{N\,s}$}
\loigiai{
Xung lượng bằng diện tích dưới đồ thị: $J=F\Delta t=24\cdot1$. Suy ra kết quả $24\ \mathrm{N\,s}$.
}
\end{ex}

% ===== Câu 118 =====
% ID: AIQ_L10C5_FULL_0118
% Mức: VDC
\begin{ex}
% ID: AIQ_L10C5_FULL_0118
% Mức: VDC
Đồ thị $F-t$ là đoạn nằm ngang ở $F=26\ \mathrm{N}$ từ $t=0$ đến $t=1,5\ \mathrm{s}$. Tính xung lượng của lực.
\choice
{$27,5\ \mathrm{N\,s}$}
{\True $39\ \mathrm{N\,s}$}
{$17,33\ \mathrm{N\,s}$}
{$19,5\ \mathrm{N\,s}$}
\loigiai{
Xung lượng bằng diện tích dưới đồ thị: $J=F\Delta t=26\cdot1,5$. Suy ra kết quả $39\ \mathrm{N\,s}$.
}
\end{ex}

% ===== Câu 119 =====
% ID: AIQ_L10C5_FULL_0119
% Mức: NB
\begin{ex}
% ID: AIQ_L10C5_FULL_0119
% Mức: NB
Xét các phát biểu thuộc dạng “Khai thác đồ thị lực – thời gian” ở tình huống số 1.
\choiceTF
{\True Diện tích dưới đồ thị $F-t$ biểu diễn xung lượng.}
{\True Đồ thị lực không đổi tạo diện tích hình chữ nhật.}
{Xung lượng không phụ thuộc thời gian tác dụng.}
{Độ dốc đồ thị $F-t$ luôn bằng động lượng.}
\loigiai{
\begin{itemchoice}
\item (Đúng) Diện tích dưới đồ thị $F-t$ biểu diễn xung lượng. (Vì): Phù hợp với định nghĩa hoặc hệ thức vật lí. b) Đúng: Phù hợp với định nghĩa hoặc hệ thức vật lí. c) Sai: Không phù hợp với định nghĩa hoặc hệ thức vật lí. d) Sai: Không phù hợp với định nghĩa hoặc hệ thức vật lí
\item (Đúng) Đồ thị lực không đổi tạo diện tích hình chữ nhật.
\item (Sai) Xung lượng không phụ thuộc thời gian tác dụng.
\item (Sai) Độ dốc đồ thị $F-t$ luôn bằng động lượng.
\end{itemchoice}
}
\end{ex}

% ===== Câu 120 =====
% ID: AIQ_L10C5_FULL_0120
% Mức: TH
\begin{ex}
% ID: AIQ_L10C5_FULL_0120
% Mức: TH
Xét các phát biểu thuộc dạng “Khai thác đồ thị lực – thời gian” ở tình huống số 2.
\choiceTF
{\True Diện tích dưới đồ thị $F-t$ biểu diễn xung lượng.}
{\True Đồ thị lực không đổi tạo diện tích hình chữ nhật.}
{Xung lượng không phụ thuộc thời gian tác dụng.}
{Độ dốc đồ thị $F-t$ luôn bằng động lượng.}
\loigiai{
\begin{itemchoice}
\item (Đúng) Diện tích dưới đồ thị $F-t$ biểu diễn xung lượng. (Vì): Phù hợp với định nghĩa hoặc hệ thức vật lí. b) Đúng: Phù hợp với định nghĩa hoặc hệ thức vật lí. c) Sai: Không phù hợp với định nghĩa hoặc hệ thức vật lí. d) Sai: Không phù hợp với định nghĩa hoặc hệ thức vật lí
\item (Đúng) Đồ thị lực không đổi tạo diện tích hình chữ nhật.
\item (Sai) Xung lượng không phụ thuộc thời gian tác dụng.
\item (Sai) Độ dốc đồ thị $F-t$ luôn bằng động lượng.
\end{itemchoice}
}
\end{ex}

% ===== Câu 121 =====
% ID: AIQ_L10C5_FULL_0121
% Mức: TH
\begin{ex}
% ID: AIQ_L10C5_FULL_0121
% Mức: TH
Xét các phát biểu thuộc dạng “Khai thác đồ thị lực – thời gian” ở tình huống số 3.
\choiceTF
{\True Diện tích dưới đồ thị $F-t$ biểu diễn xung lượng.}
{\True Đồ thị lực không đổi tạo diện tích hình chữ nhật.}
{Xung lượng không phụ thuộc thời gian tác dụng.}
{Độ dốc đồ thị $F-t$ luôn bằng động lượng.}
\loigiai{
\begin{itemchoice}
\item (Đúng) Diện tích dưới đồ thị $F-t$ biểu diễn xung lượng. (Vì): Phù hợp với định nghĩa hoặc hệ thức vật lí. b) Đúng: Phù hợp với định nghĩa hoặc hệ thức vật lí. c) Sai: Không phù hợp với định nghĩa hoặc hệ thức vật lí. d) Sai: Không phù hợp với định nghĩa hoặc hệ thức vật lí
\item (Đúng) Đồ thị lực không đổi tạo diện tích hình chữ nhật.
\item (Sai) Xung lượng không phụ thuộc thời gian tác dụng.
\item (Sai) Độ dốc đồ thị $F-t$ luôn bằng động lượng.
\end{itemchoice}
}
\end{ex}

% ===== Câu 122 =====
% ID: AIQ_L10C5_FULL_0122
% Mức: VD
\begin{ex}
% ID: AIQ_L10C5_FULL_0122
% Mức: VD
Xét các phát biểu thuộc dạng “Khai thác đồ thị lực – thời gian” ở tình huống số 4.
\choiceTF
{\True Diện tích dưới đồ thị $F-t$ biểu diễn xung lượng.}
{\True Đồ thị lực không đổi tạo diện tích hình chữ nhật.}
{Xung lượng không phụ thuộc thời gian tác dụng.}
{Độ dốc đồ thị $F-t$ luôn bằng động lượng.}
\loigiai{
\begin{itemchoice}
\item (Đúng) Diện tích dưới đồ thị $F-t$ biểu diễn xung lượng. (Vì): Phù hợp với định nghĩa hoặc hệ thức vật lí. b) Đúng: Phù hợp với định nghĩa hoặc hệ thức vật lí. c) Sai: Không phù hợp với định nghĩa hoặc hệ thức vật lí. d) Sai: Không phù hợp với định nghĩa hoặc hệ thức vật lí
\item (Đúng) Đồ thị lực không đổi tạo diện tích hình chữ nhật.
\item (Sai) Xung lượng không phụ thuộc thời gian tác dụng.
\item (Sai) Độ dốc đồ thị $F-t$ luôn bằng động lượng.
\end{itemchoice}
}
\end{ex}

% ===== Câu 123 =====
% ID: AIQ_L10C5_FULL_0123
% Mức: VDC
\begin{ex}
% ID: AIQ_L10C5_FULL_0123
% Mức: VDC
Xét các phát biểu thuộc dạng “Khai thác đồ thị lực – thời gian” ở tình huống số 5.
\choiceTF
{\True Diện tích dưới đồ thị $F-t$ biểu diễn xung lượng.}
{\True Đồ thị lực không đổi tạo diện tích hình chữ nhật.}
{Xung lượng không phụ thuộc thời gian tác dụng.}
{Độ dốc đồ thị $F-t$ luôn bằng động lượng.}
\loigiai{
\begin{itemchoice}
\item (Đúng) Diện tích dưới đồ thị $F-t$ biểu diễn xung lượng. (Vì): Phù hợp với định nghĩa hoặc hệ thức vật lí. b) Đúng: Phù hợp với định nghĩa hoặc hệ thức vật lí. c) Sai: Không phù hợp với định nghĩa hoặc hệ thức vật lí. d) Sai: Không phù hợp với định nghĩa hoặc hệ thức vật lí
\item (Đúng) Đồ thị lực không đổi tạo diện tích hình chữ nhật.
\item (Sai) Xung lượng không phụ thuộc thời gian tác dụng.
\item (Sai) Độ dốc đồ thị $F-t$ luôn bằng động lượng.
\end{itemchoice}
}
\end{ex}

% ===== Câu 124 =====
% ID: AIQ_L10C5_FULL_0124
% Mức: TH
\begin{ex}
% ID: AIQ_L10C5_FULL_0124
% Mức: TH
Đồ thị $F-t$ là đoạn nằm ngang ở $F=28\ \mathrm{N}$ từ $t=0$ đến $t=2\ \mathrm{s}$. Tính xung lượng của lực.
\shortans{56}
\loigiai{
Xung lượng bằng diện tích dưới đồ thị: $J=F\Delta t=28\cdot2$. Suy ra kết quả $56\ \mathrm{N\,s}$. Đáp án trả lời ngắn không ghi đơn vị.
}
\end{ex}

% ===== Câu 125 =====
% ID: AIQ_L10C5_FULL_0125
% Mức: TH
\begin{ex}
% ID: AIQ_L10C5_FULL_0125
% Mức: TH
Đồ thị $F-t$ là đoạn nằm ngang ở $F=30\ \mathrm{N}$ từ $t=0$ đến $t=2,5\ \mathrm{s}$. Tính xung lượng của lực.
\shortans{75}
\loigiai{
Xung lượng bằng diện tích dưới đồ thị: $J=F\Delta t=30\cdot2,5$. Suy ra kết quả $75\ \mathrm{N\,s}$. Đáp án trả lời ngắn không ghi đơn vị.
}
\end{ex}

% ===== Câu 126 =====
% ID: AIQ_L10C5_FULL_0126
% Mức: VD
\begin{ex}
% ID: AIQ_L10C5_FULL_0126
% Mức: VD
Đồ thị $F-t$ là đoạn nằm ngang ở $F=32\ \mathrm{N}$ từ $t=0$ đến $t=1\ \mathrm{s}$. Tính xung lượng của lực.
\shortans{32}
\loigiai{
Xung lượng bằng diện tích dưới đồ thị: $J=F\Delta t=32\cdot1$. Suy ra kết quả $32\ \mathrm{N\,s}$. Đáp án trả lời ngắn không ghi đơn vị.
}
\end{ex}

% ===== Câu 127 =====
% ID: AIQ_L10C5_FULL_0127
% Mức: VD
\begin{ex}
% ID: AIQ_L10C5_FULL_0127
% Mức: VD
Đồ thị $F-t$ là đoạn nằm ngang ở $F=34\ \mathrm{N}$ từ $t=0$ đến $t=1,5\ \mathrm{s}$. Tính xung lượng của lực.
\shortans{51}
\loigiai{
Xung lượng bằng diện tích dưới đồ thị: $J=F\Delta t=34\cdot1,5$. Suy ra kết quả $51\ \mathrm{N\,s}$. Đáp án trả lời ngắn không ghi đơn vị.
}
\end{ex}

% ===== Câu 128 =====
% ID: AIQ_L10C5_FULL_0128
% Mức: VDC
\begin{ex}
% ID: AIQ_L10C5_FULL_0128
% Mức: VDC
Đồ thị $F-t$ là đoạn nằm ngang ở $F=36\ \mathrm{N}$ từ $t=0$ đến $t=2\ \mathrm{s}$. Tính xung lượng của lực.
\shortans{72}
\loigiai{
Xung lượng bằng diện tích dưới đồ thị: $J=F\Delta t=36\cdot2$. Suy ra kết quả $72\ \mathrm{N\,s}$. Đáp án trả lời ngắn không ghi đơn vị.
}
\end{ex}

% ===== Câu 129 =====
% ID: AIQ_L10C5_FULL_0129
% Mức: VDC
\begin{ex}
% ID: AIQ_L10C5_FULL_0129
% Mức: VDC
Đồ thị $F-t$ là đoạn nằm ngang ở $F=38\ \mathrm{N}$ từ $t=0$ đến $t=2,5\ \mathrm{s}$. Tính xung lượng của lực.
\shortans{95}
\loigiai{
Xung lượng bằng diện tích dưới đồ thị: $J=F\Delta t=38\cdot2,5$. Suy ra kết quả $95\ \mathrm{N\,s}$. Đáp án trả lời ngắn không ghi đơn vị.
}
\end{ex}

% ===== Câu 130 =====
% ID: AIQ_L10C5_FULL_0130
% Mức: TH
\begin{bt}
% ID: AIQ_L10C5_FULL_0130
% Mức: TH
Trình bày cách giải: Đồ thị $F-t$ là đoạn nằm ngang ở $F=40\ \mathrm{N}$ từ $t=0$ đến $t=1\ \mathrm{s}$. Tính xung lượng của lực.
\loigiai{
Xung lượng bằng diện tích dưới đồ thị: $J=F\Delta t=40\cdot1$. Suy ra kết quả $40\ \mathrm{N\,s}$.
}
\end{bt}

% ===== Câu 131 =====
% ID: AIQ_L10C5_FULL_0131
% Mức: TH
\begin{bt}
% ID: AIQ_L10C5_FULL_0131
% Mức: TH
Trình bày cách giải: Đồ thị $F-t$ là đoạn nằm ngang ở $F=42\ \mathrm{N}$ từ $t=0$ đến $t=1,5\ \mathrm{s}$. Tính xung lượng của lực.
\loigiai{
Xung lượng bằng diện tích dưới đồ thị: $J=F\Delta t=42\cdot1,5$. Suy ra kết quả $63\ \mathrm{N\,s}$.
}
\end{bt}

% ===== Câu 132 =====
% ID: AIQ_L10C5_FULL_0132
% Mức: VD
\begin{bt}
% ID: AIQ_L10C5_FULL_0132
% Mức: VD
Trình bày cách giải: Đồ thị $F-t$ là đoạn nằm ngang ở $F=44\ \mathrm{N}$ từ $t=0$ đến $t=2\ \mathrm{s}$. Tính xung lượng của lực.
\loigiai{
Xung lượng bằng diện tích dưới đồ thị: $J=F\Delta t=44\cdot2$. Suy ra kết quả $88\ \mathrm{N\,s}$.
}
\end{bt}

% ===== Câu 133 =====
% ID: AIQ_L10C5_FULL_0133
% Mức: VD
\begin{bt}
% ID: AIQ_L10C5_FULL_0133
% Mức: VD
Trình bày cách giải: Đồ thị $F-t$ là đoạn nằm ngang ở $F=46\ \mathrm{N}$ từ $t=0$ đến $t=2,5\ \mathrm{s}$. Tính xung lượng của lực.
\loigiai{
Xung lượng bằng diện tích dưới đồ thị: $J=F\Delta t=46\cdot2,5$. Suy ra kết quả $115\ \mathrm{N\,s}$.
}
\end{bt}

% ===== Câu 134 =====
% ID: AIQ_L10C5_FULL_0134
% Mức: VDC
\begin{bt}
% ID: AIQ_L10C5_FULL_0134
% Mức: VDC
Trình bày cách giải: Đồ thị $F-t$ là đoạn nằm ngang ở $F=48\ \mathrm{N}$ từ $t=0$ đến $t=1\ \mathrm{s}$. Tính xung lượng của lực.
\loigiai{
Xung lượng bằng diện tích dưới đồ thị: $J=F\Delta t=48\cdot1$. Suy ra kết quả $48\ \mathrm{N\,s}$.
}
\end{bt}

% ===== Câu 135 =====
% ID: AIQ_L10C5_FULL_0135
% Mức: VDC
\begin{bt}
% ID: AIQ_L10C5_FULL_0135
% Mức: VDC
Trình bày cách giải: Đồ thị $F-t$ là đoạn nằm ngang ở $F=50\ \mathrm{N}$ từ $t=0$ đến $t=1,5\ \mathrm{s}$. Tính xung lượng của lực.
\loigiai{
Xung lượng bằng diện tích dưới đồ thị: $J=F\Delta t=50\cdot1,5$. Suy ra kết quả $75\ \mathrm{N\,s}$.
}
\end{bt}

% ===== Câu 136 =====
% ID: AIQ_L10C5_FULL_0136
% Mức: NB
\dangbt{Bài toán động lượng và xung lượng gắn với thực tiễn}
\begin{ex}
% ID: AIQ_L10C5_FULL_0136
% Mức: NB
Người khối lượng $50\ \mathrm{kg}$ đang chuyển động với tốc độ $4\ \mathrm{m/s}$ được túi khí làm dừng trong $0,5\ \mathrm{s}$. Bỏ qua lực khác, tính độ lớn lực trung bình.
\choice
{\True $400\ \mathrm{N}$}
{$200\ \mathrm{N}$}
{$100\ \mathrm{N}$}
{$100\ \mathrm{N}$}
\loigiai{
$F_{tb}=|\Delta p|/\Delta t=mv/\Delta t=50\cdot4/0,5$. Suy ra kết quả $400\ \mathrm{N}$.
}
\end{ex}

% ===== Câu 137 =====
% ID: AIQ_L10C5_FULL_0137
% Mức: NB
\begin{ex}
% ID: AIQ_L10C5_FULL_0137
% Mức: NB
Người khối lượng $55\ \mathrm{kg}$ đang chuyển động với tốc độ $5\ \mathrm{m/s}$ được túi khí làm dừng trong $0,6\ \mathrm{s}$. Bỏ qua lực khác, tính độ lớn lực trung bình.
\choice
{$275\ \mathrm{N}$}
{\True $458,33\ \mathrm{N}$}
{$165\ \mathrm{N}$}
{$91,67\ \mathrm{N}$}
\loigiai{
$F_{tb}=|\Delta p|/\Delta t=mv/\Delta t=55\cdot5/0,6$. Suy ra kết quả $458,33\ \mathrm{N}$.
}
\end{ex}

% ===== Câu 138 =====
% ID: AIQ_L10C5_FULL_0138
% Mức: TH
\begin{ex}
% ID: AIQ_L10C5_FULL_0138
% Mức: TH
Người khối lượng $60\ \mathrm{kg}$ đang chuyển động với tốc độ $6\ \mathrm{m/s}$ được túi khí làm dừng trong $0,7\ \mathrm{s}$. Bỏ qua lực khác, tính độ lớn lực trung bình.
\choice
{$360\ \mathrm{N}$}
{$252\ \mathrm{N}$}
{\True $514,29\ \mathrm{N}$}
{$85,71\ \mathrm{N}$}
\loigiai{
$F_{tb}=|\Delta p|/\Delta t=mv/\Delta t=60\cdot6/0,7$. Suy ra kết quả $514,29\ \mathrm{N}$.
}
\end{ex}

% ===== Câu 139 =====
% ID: AIQ_L10C5_FULL_0139
% Mức: TH
\begin{ex}
% ID: AIQ_L10C5_FULL_0139
% Mức: TH
Người khối lượng $65\ \mathrm{kg}$ đang chuyển động với tốc độ $7\ \mathrm{m/s}$ được túi khí làm dừng trong $0,8\ \mathrm{s}$. Bỏ qua lực khác, tính độ lớn lực trung bình.
\choice
{$455\ \mathrm{N}$}
{$364\ \mathrm{N}$}
{$81,25\ \mathrm{N}$}
{\True $568,75\ \mathrm{N}$}
\loigiai{
$F_{tb}=|\Delta p|/\Delta t=mv/\Delta t=65\cdot7/0,8$. Suy ra kết quả $568,75\ \mathrm{N}$.
}
\end{ex}

% ===== Câu 140 =====
% ID: AIQ_L10C5_FULL_0140
% Mức: VD
\begin{ex}
% ID: AIQ_L10C5_FULL_0140
% Mức: VD
Người khối lượng $70\ \mathrm{kg}$ đang chuyển động với tốc độ $4\ \mathrm{m/s}$ được túi khí làm dừng trong $0,9\ \mathrm{s}$. Bỏ qua lực khác, tính độ lớn lực trung bình.
\choice
{\True $311,11\ \mathrm{N}$}
{$280\ \mathrm{N}$}
{$252\ \mathrm{N}$}
{$77,78\ \mathrm{N}$}
\loigiai{
$F_{tb}=|\Delta p|/\Delta t=mv/\Delta t=70\cdot4/0,9$. Suy ra kết quả $311,11\ \mathrm{N}$.
}
\end{ex}

% ===== Câu 141 =====
% ID: AIQ_L10C5_FULL_0141
% Mức: VD
\begin{ex}
% ID: AIQ_L10C5_FULL_0141
% Mức: VD
Người khối lượng $75\ \mathrm{kg}$ đang chuyển động với tốc độ $5\ \mathrm{m/s}$ được túi khí làm dừng trong $0,5\ \mathrm{s}$. Bỏ qua lực khác, tính độ lớn lực trung bình.
\choice
{$375\ \mathrm{N}$}
{\True $750\ \mathrm{N}$}
{$187,5\ \mathrm{N}$}
{$150\ \mathrm{N}$}
\loigiai{
$F_{tb}=|\Delta p|/\Delta t=mv/\Delta t=75\cdot5/0,5$. Suy ra kết quả $750\ \mathrm{N}$.
}
\end{ex}

% ===== Câu 142 =====
% ID: AIQ_L10C5_FULL_0142
% Mức: VD
\begin{ex}
% ID: AIQ_L10C5_FULL_0142
% Mức: VD
Người khối lượng $80\ \mathrm{kg}$ đang chuyển động với tốc độ $6\ \mathrm{m/s}$ được túi khí làm dừng trong $0,6\ \mathrm{s}$. Bỏ qua lực khác, tính độ lớn lực trung bình.
\choice
{$480\ \mathrm{N}$}
{$288\ \mathrm{N}$}
{\True $800\ \mathrm{N}$}
{$133,33\ \mathrm{N}$}
\loigiai{
$F_{tb}=|\Delta p|/\Delta t=mv/\Delta t=80\cdot6/0,6$. Suy ra kết quả $800\ \mathrm{N}$.
}
\end{ex}

% ===== Câu 143 =====
% ID: AIQ_L10C5_FULL_0143
% Mức: VDC
\begin{ex}
% ID: AIQ_L10C5_FULL_0143
% Mức: VDC
Người khối lượng $85\ \mathrm{kg}$ đang chuyển động với tốc độ $7\ \mathrm{m/s}$ được túi khí làm dừng trong $0,7\ \mathrm{s}$. Bỏ qua lực khác, tính độ lớn lực trung bình.
\choice
{$595\ \mathrm{N}$}
{$416,5\ \mathrm{N}$}
{$121,43\ \mathrm{N}$}
{\True $850\ \mathrm{N}$}
\loigiai{
$F_{tb}=|\Delta p|/\Delta t=mv/\Delta t=85\cdot7/0,7$. Suy ra kết quả $850\ \mathrm{N}$.
}
\end{ex}

% ===== Câu 144 =====
% ID: AIQ_L10C5_FULL_0144
% Mức: VDC
\begin{ex}
% ID: AIQ_L10C5_FULL_0144
% Mức: VDC
Người khối lượng $90\ \mathrm{kg}$ đang chuyển động với tốc độ $4\ \mathrm{m/s}$ được túi khí làm dừng trong $0,8\ \mathrm{s}$. Bỏ qua lực khác, tính độ lớn lực trung bình.
\choice
{\True $450\ \mathrm{N}$}
{$360\ \mathrm{N}$}
{$288\ \mathrm{N}$}
{$112,5\ \mathrm{N}$}
\loigiai{
$F_{tb}=|\Delta p|/\Delta t=mv/\Delta t=90\cdot4/0,8$. Suy ra kết quả $450\ \mathrm{N}$.
}
\end{ex}

% ===== Câu 145 =====
% ID: AIQ_L10C5_FULL_0145
% Mức: VDC
\begin{ex}
% ID: AIQ_L10C5_FULL_0145
% Mức: VDC
Người khối lượng $95\ \mathrm{kg}$ đang chuyển động với tốc độ $5\ \mathrm{m/s}$ được túi khí làm dừng trong $0,9\ \mathrm{s}$. Bỏ qua lực khác, tính độ lớn lực trung bình.
\choice
{$475\ \mathrm{N}$}
{\True $527,78\ \mathrm{N}$}
{$427,5\ \mathrm{N}$}
{$105,56\ \mathrm{N}$}
\loigiai{
$F_{tb}=|\Delta p|/\Delta t=mv/\Delta t=95\cdot5/0,9$. Suy ra kết quả $527,78\ \mathrm{N}$.
}
\end{ex}

% ===== Câu 146 =====
% ID: AIQ_L10C5_FULL_0146
% Mức: NB
\begin{ex}
% ID: AIQ_L10C5_FULL_0146
% Mức: NB
Xét các phát biểu thuộc dạng “Bài toán động lượng và xung lượng gắn với thực tiễn” ở tình huống số 1.
\choiceTF
{\True Túi khí làm tăng thời gian biến đổi động lượng.}
{\True Tăng thời gian dừng giúp giảm lực trung bình.}
{Đệm hơi làm người dừng tức thời.}
{\True Găng tay bảo hộ có thể giảm lực va chạm trung bình.}
\loigiai{
\begin{itemchoice}
\item (Đúng) Túi khí làm tăng thời gian biến đổi động lượng. (Vì): Phù hợp với định nghĩa hoặc hệ thức vật lí. b) Đúng: Phù hợp với định nghĩa hoặc hệ thức vật lí. c) Sai: Không phù hợp với định nghĩa hoặc hệ thức vật lí. d) Đúng: Phù hợp với định nghĩa hoặc hệ thức vật lí
\item (Đúng) Tăng thời gian dừng giúp giảm lực trung bình.
\item (Sai) Đệm hơi làm người dừng tức thời.
\item (Đúng) Găng tay bảo hộ có thể giảm lực va chạm trung bình.
\end{itemchoice}
}
\end{ex}

% ===== Câu 147 =====
% ID: AIQ_L10C5_FULL_0147
% Mức: TH
\begin{ex}
% ID: AIQ_L10C5_FULL_0147
% Mức: TH
Xét các phát biểu thuộc dạng “Bài toán động lượng và xung lượng gắn với thực tiễn” ở tình huống số 2.
\choiceTF
{\True Túi khí làm tăng thời gian biến đổi động lượng.}
{\True Tăng thời gian dừng giúp giảm lực trung bình.}
{Đệm hơi làm người dừng tức thời.}
{\True Găng tay bảo hộ có thể giảm lực va chạm trung bình.}
\loigiai{
\begin{itemchoice}
\item (Đúng) Túi khí làm tăng thời gian biến đổi động lượng. (Vì): Phù hợp với định nghĩa hoặc hệ thức vật lí. b) Đúng: Phù hợp với định nghĩa hoặc hệ thức vật lí. c) Sai: Không phù hợp với định nghĩa hoặc hệ thức vật lí. d) Đúng: Phù hợp với định nghĩa hoặc hệ thức vật lí
\item (Đúng) Tăng thời gian dừng giúp giảm lực trung bình.
\item (Sai) Đệm hơi làm người dừng tức thời.
\item (Đúng) Găng tay bảo hộ có thể giảm lực va chạm trung bình.
\end{itemchoice}
}
\end{ex}

% ===== Câu 148 =====
% ID: AIQ_L10C5_FULL_0148
% Mức: TH
\begin{ex}
% ID: AIQ_L10C5_FULL_0148
% Mức: TH
Xét các phát biểu thuộc dạng “Bài toán động lượng và xung lượng gắn với thực tiễn” ở tình huống số 3.
\choiceTF
{\True Túi khí làm tăng thời gian biến đổi động lượng.}
{\True Tăng thời gian dừng giúp giảm lực trung bình.}
{Đệm hơi làm người dừng tức thời.}
{\True Găng tay bảo hộ có thể giảm lực va chạm trung bình.}
\loigiai{
\begin{itemchoice}
\item (Đúng) Túi khí làm tăng thời gian biến đổi động lượng. (Vì): Phù hợp với định nghĩa hoặc hệ thức vật lí. b) Đúng: Phù hợp với định nghĩa hoặc hệ thức vật lí. c) Sai: Không phù hợp với định nghĩa hoặc hệ thức vật lí. d) Đúng: Phù hợp với định nghĩa hoặc hệ thức vật lí
\item (Đúng) Tăng thời gian dừng giúp giảm lực trung bình.
\item (Sai) Đệm hơi làm người dừng tức thời.
\item (Đúng) Găng tay bảo hộ có thể giảm lực va chạm trung bình.
\end{itemchoice}
}
\end{ex}

% ===== Câu 149 =====
% ID: AIQ_L10C5_FULL_0149
% Mức: VD
\begin{ex}
% ID: AIQ_L10C5_FULL_0149
% Mức: VD
Xét các phát biểu thuộc dạng “Bài toán động lượng và xung lượng gắn với thực tiễn” ở tình huống số 4.
\choiceTF
{\True Túi khí làm tăng thời gian biến đổi động lượng.}
{\True Tăng thời gian dừng giúp giảm lực trung bình.}
{Đệm hơi làm người dừng tức thời.}
{\True Găng tay bảo hộ có thể giảm lực va chạm trung bình.}
\loigiai{
\begin{itemchoice}
\item (Đúng) Túi khí làm tăng thời gian biến đổi động lượng. (Vì): Phù hợp với định nghĩa hoặc hệ thức vật lí. b) Đúng: Phù hợp với định nghĩa hoặc hệ thức vật lí. c) Sai: Không phù hợp với định nghĩa hoặc hệ thức vật lí. d) Đúng: Phù hợp với định nghĩa hoặc hệ thức vật lí
\item (Đúng) Tăng thời gian dừng giúp giảm lực trung bình.
\item (Sai) Đệm hơi làm người dừng tức thời.
\item (Đúng) Găng tay bảo hộ có thể giảm lực va chạm trung bình.
\end{itemchoice}
}
\end{ex}

% ===== Câu 150 =====
% ID: AIQ_L10C5_FULL_0150
% Mức: VDC
\begin{ex}
% ID: AIQ_L10C5_FULL_0150
% Mức: VDC
Xét các phát biểu thuộc dạng “Bài toán động lượng và xung lượng gắn với thực tiễn” ở tình huống số 5.
\choiceTF
{\True Túi khí làm tăng thời gian biến đổi động lượng.}
{\True Tăng thời gian dừng giúp giảm lực trung bình.}
{Đệm hơi làm người dừng tức thời.}
{\True Găng tay bảo hộ có thể giảm lực va chạm trung bình.}
\loigiai{
\begin{itemchoice}
\item (Đúng) Túi khí làm tăng thời gian biến đổi động lượng. (Vì): Phù hợp với định nghĩa hoặc hệ thức vật lí. b) Đúng: Phù hợp với định nghĩa hoặc hệ thức vật lí. c) Sai: Không phù hợp với định nghĩa hoặc hệ thức vật lí. d) Đúng: Phù hợp với định nghĩa hoặc hệ thức vật lí
\item (Đúng) Tăng thời gian dừng giúp giảm lực trung bình.
\item (Sai) Đệm hơi làm người dừng tức thời.
\item (Đúng) Găng tay bảo hộ có thể giảm lực va chạm trung bình.
\end{itemchoice}
}
\end{ex}

% ===== Câu 151 =====
% ID: AIQ_L10C5_FULL_0151
% Mức: TH
\begin{ex}
% ID: AIQ_L10C5_FULL_0151
% Mức: TH
Người khối lượng $100\ \mathrm{kg}$ đang chuyển động với tốc độ $6\ \mathrm{m/s}$ được túi khí làm dừng trong $0,5\ \mathrm{s}$. Bỏ qua lực khác, tính độ lớn lực trung bình.
\shortans{1200}
\loigiai{
$F_{tb}=|\Delta p|/\Delta t=mv/\Delta t=100\cdot6/0,5$. Suy ra kết quả $1200\ \mathrm{N}$. Đáp án trả lời ngắn không ghi đơn vị.
}
\end{ex}

% ===== Câu 152 =====
% ID: AIQ_L10C5_FULL_0152
% Mức: TH
\begin{ex}
% ID: AIQ_L10C5_FULL_0152
% Mức: TH
Người khối lượng $105\ \mathrm{kg}$ đang chuyển động với tốc độ $7\ \mathrm{m/s}$ được túi khí làm dừng trong $0,6\ \mathrm{s}$. Bỏ qua lực khác, tính độ lớn lực trung bình.
\shortans{1225}
\loigiai{
$F_{tb}=|\Delta p|/\Delta t=mv/\Delta t=105\cdot7/0,6$. Suy ra kết quả $1225\ \mathrm{N}$. Đáp án trả lời ngắn không ghi đơn vị.
}
\end{ex}

% ===== Câu 153 =====
% ID: AIQ_L10C5_FULL_0153
% Mức: VD
\begin{ex}
% ID: AIQ_L10C5_FULL_0153
% Mức: VD
Người khối lượng $110\ \mathrm{kg}$ đang chuyển động với tốc độ $4\ \mathrm{m/s}$ được túi khí làm dừng trong $0,7\ \mathrm{s}$. Bỏ qua lực khác, tính độ lớn lực trung bình.
\shortans{628,57}
\loigiai{
$F_{tb}=|\Delta p|/\Delta t=mv/\Delta t=110\cdot4/0,7$. Suy ra kết quả $628,57\ \mathrm{N}$. Đáp án trả lời ngắn không ghi đơn vị.
}
\end{ex}

% ===== Câu 154 =====
% ID: AIQ_L10C5_FULL_0154
% Mức: VD
\begin{ex}
% ID: AIQ_L10C5_FULL_0154
% Mức: VD
Người khối lượng $115\ \mathrm{kg}$ đang chuyển động với tốc độ $5\ \mathrm{m/s}$ được túi khí làm dừng trong $0,8\ \mathrm{s}$. Bỏ qua lực khác, tính độ lớn lực trung bình.
\shortans{718,75}
\loigiai{
$F_{tb}=|\Delta p|/\Delta t=mv/\Delta t=115\cdot5/0,8$. Suy ra kết quả $718,75\ \mathrm{N}$. Đáp án trả lời ngắn không ghi đơn vị.
}
\end{ex}

% ===== Câu 155 =====
% ID: AIQ_L10C5_FULL_0155
% Mức: VDC
\begin{ex}
% ID: AIQ_L10C5_FULL_0155
% Mức: VDC
Người khối lượng $120\ \mathrm{kg}$ đang chuyển động với tốc độ $6\ \mathrm{m/s}$ được túi khí làm dừng trong $0,9\ \mathrm{s}$. Bỏ qua lực khác, tính độ lớn lực trung bình.
\shortans{800}
\loigiai{
$F_{tb}=|\Delta p|/\Delta t=mv/\Delta t=120\cdot6/0,9$. Suy ra kết quả $800\ \mathrm{N}$. Đáp án trả lời ngắn không ghi đơn vị.
}
\end{ex}

% ===== Câu 156 =====
% ID: AIQ_L10C5_FULL_0156
% Mức: VDC
\begin{ex}
% ID: AIQ_L10C5_FULL_0156
% Mức: VDC
Người khối lượng $125\ \mathrm{kg}$ đang chuyển động với tốc độ $7\ \mathrm{m/s}$ được túi khí làm dừng trong $0,5\ \mathrm{s}$. Bỏ qua lực khác, tính độ lớn lực trung bình.
\shortans{1750}
\loigiai{
$F_{tb}=|\Delta p|/\Delta t=mv/\Delta t=125\cdot7/0,5$. Suy ra kết quả $1750\ \mathrm{N}$. Đáp án trả lời ngắn không ghi đơn vị.
}
\end{ex}

% ===== Câu 157 =====
% ID: AIQ_L10C5_FULL_0157
% Mức: TH
\begin{bt}
% ID: AIQ_L10C5_FULL_0157
% Mức: TH
Trình bày cách giải: Người khối lượng $130\ \mathrm{kg}$ đang chuyển động với tốc độ $4\ \mathrm{m/s}$ được túi khí làm dừng trong $0,6\ \mathrm{s}$. Bỏ qua lực khác, tính độ lớn lực trung bình.
\loigiai{
$F_{tb}=|\Delta p|/\Delta t=mv/\Delta t=130\cdot4/0,6$. Suy ra kết quả $866,67\ \mathrm{N}$.
}
\end{bt}

% ===== Câu 158 =====
% ID: AIQ_L10C5_FULL_0158
% Mức: TH
\begin{bt}
% ID: AIQ_L10C5_FULL_0158
% Mức: TH
Trình bày cách giải: Người khối lượng $135\ \mathrm{kg}$ đang chuyển động với tốc độ $5\ \mathrm{m/s}$ được túi khí làm dừng trong $0,7\ \mathrm{s}$. Bỏ qua lực khác, tính độ lớn lực trung bình.
\loigiai{
$F_{tb}=|\Delta p|/\Delta t=mv/\Delta t=135\cdot5/0,7$. Suy ra kết quả $964,29\ \mathrm{N}$.
}
\end{bt}

% ===== Câu 159 =====
% ID: AIQ_L10C5_FULL_0159
% Mức: VD
\begin{bt}
% ID: AIQ_L10C5_FULL_0159
% Mức: VD
Trình bày cách giải: Người khối lượng $140\ \mathrm{kg}$ đang chuyển động với tốc độ $6\ \mathrm{m/s}$ được túi khí làm dừng trong $0,8\ \mathrm{s}$. Bỏ qua lực khác, tính độ lớn lực trung bình.
\loigiai{
$F_{tb}=|\Delta p|/\Delta t=mv/\Delta t=140\cdot6/0,8$. Suy ra kết quả $1050\ \mathrm{N}$.
}
\end{bt}

% ===== Câu 160 =====
% ID: AIQ_L10C5_FULL_0160
% Mức: VD
\begin{bt}
% ID: AIQ_L10C5_FULL_0160
% Mức: VD
Trình bày cách giải: Người khối lượng $145\ \mathrm{kg}$ đang chuyển động với tốc độ $7\ \mathrm{m/s}$ được túi khí làm dừng trong $0,9\ \mathrm{s}$. Bỏ qua lực khác, tính độ lớn lực trung bình.
\loigiai{
$F_{tb}=|\Delta p|/\Delta t=mv/\Delta t=145\cdot7/0,9$. Suy ra kết quả $1127,78\ \mathrm{N}$.
}
\end{bt}

% ===== Câu 161 =====
% ID: AIQ_L10C5_FULL_0161
% Mức: VDC
\begin{bt}
% ID: AIQ_L10C5_FULL_0161
% Mức: VDC
Trình bày cách giải: Người khối lượng $150\ \mathrm{kg}$ đang chuyển động với tốc độ $4\ \mathrm{m/s}$ được túi khí làm dừng trong $0,5\ \mathrm{s}$. Bỏ qua lực khác, tính độ lớn lực trung bình.
\loigiai{
$F_{tb}=|\Delta p|/\Delta t=mv/\Delta t=150\cdot4/0,5$. Suy ra kết quả $1200\ \mathrm{N}$.
}
\end{bt}

% ===== Câu 162 =====
% ID: AIQ_L10C5_FULL_0162
% Mức: VDC
\begin{bt}
% ID: AIQ_L10C5_FULL_0162
% Mức: VDC
Trình bày cách giải: Người khối lượng $155\ \mathrm{kg}$ đang chuyển động với tốc độ $5\ \mathrm{m/s}$ được túi khí làm dừng trong $0,6\ \mathrm{s}$. Bỏ qua lực khác, tính độ lớn lực trung bình.
\loigiai{
$F_{tb}=|\Delta p|/\Delta t=mv/\Delta t=155\cdot5/0,6$. Suy ra kết quả $1291,67\ \mathrm{N}$.
}
\end{bt}
